% !TeX program = xelatex
\documentclass[11pt,a4paper,fontset=fandol]{ctexart}
\usepackage[margin=18mm]{geometry}
\usepackage{amsmath,amssymb,bm,enumitem,hyperref}
\hypersetup{hidelinks,pdftitle={大一数学分析期末复习·下册}}
\setlist[enumerate]{label=（\arabic*）,leftmargin=2.6em,itemsep=0.25em}
\newcommand{\dd}{\,\mathrm{d}}
\newenvironment{exercise}[2]{%
  \par\medskip\noindent\textbf{第#1题}\quad\textit{#2}\par\smallskip
}{\par\medskip}
\allowdisplaybreaks
\setcounter{secnumdepth}{-1}
\sloppy
\setlength{\emergencystretch}{4em}
\title{大一数学分析期末复习·下册\\[0.4em]\large 基础计算与常见题型}
\author{}
\date{2026年10月20日}
\begin{document}
\maketitle
\section*{使用说明}
本册用于大一期末分章复习，覆盖向量与空间几何、多元微分、重积分、曲线积分、曲面积分及无穷级数。题目以基础计算、常见应用和课堂短证明为主，证明题使用已学定理并说明关键条件。

全册共818题，按阅读顺序统一编号1至818。34节各保留两道补充练习，用于基础变式与常规小综合，均纳入连续编号。本版设“基础”“常规”两档：“基础”表示直接运用概念、公式或标准方法，“常规”表示期末常见的两步计算、应用或短证明。建议先做基础题，再做常规题。

常规版与iPad横屏版共用同一份题面，后者按每页一道题设计，下方纯空白用于草稿。本次仅交付LaTeX源文件，按要求未编译，未附答案。未获得本校期末试卷，难度按通用的大一期末复习要求调整；具体范围以课堂教学为准，傅里叶级数等内容按课程实际教学选用。
\tableofcontents
\clearpage
\section{第8章 向量代数与空间解析几何}
\subsection{向量及其线性运算}
\begin{exercise}{1}{基础·向量线性组合}
设 $\vec u=\vec a-\vec b+\vec c$、$\vec v=\vec a+3\vec b-2\vec c$。用 $\vec a,\vec b,\vec c$ 表示 $2\vec u-\vec v$。
\end{exercise}
\begin{exercise}{2}{基础·坐标运算}
设 $\vec a=(0,-1,1)$、$\vec b=(1,2,0)$、$\vec c=(2,0,1)$。求 $\vec a-\vec b+\vec c$ 及其模。
\end{exercise}
\begin{exercise}{3}{基础·三等分点}
已知 $A(0,0,0)$、$B(6,3,-3)$，求线段 $AB$ 的两个三等分点及线段长度。
\end{exercise}
\begin{exercise}{4}{基础·线段内分点}
已知 $A(1,0,2)$、$B(4,6,-1)$，点 $M$ 在线段 $AB$ 上且 $|AM|:|MB|=1:2$。求 $M$ 的坐标。
\end{exercise}
\begin{exercise}{5}{基础·向量平行条件}
设 $A(1,2,3)$、$B(3,5,t)$，$\vec a=(2,3,-2)$。求使 $\overrightarrow{AB}$ 与 $\vec a$ 平行的 $t$。
\end{exercise}
\begin{exercise}{6}{基础·平行向量坐标}
设 $\vec a=(1,2,3)$、$\vec b=(2,k,t)$。求使两向量平行的 $k,t$。
\end{exercise}
\begin{exercise}{7}{基础·单位向量}
设 $A(-1,2,2)$、$B(1,2,0)$。求与 $\overrightarrow{AB}$ 反向的单位向量。
\end{exercise}
\begin{exercise}{8}{基础·方向余弦}
设 $\vec a=(1,1,\sqrt2)$，求其模及三个方向余弦。
\end{exercise}
\begin{exercise}{9}{基础·标量投影}
向量 $\vec a$ 的模为 $6$，与某有向轴的夹角为 $\pi/3$。求它在该轴上的标量投影。
\end{exercise}
\begin{exercise}{10}{常规·模与正交条件}
设 $\vec a=(1,2,3)$、$\vec b=(2,-1,p)$。由 $|\vec a+\vec b|=|\vec a-\vec b|$ 求 $p$。提示：将两边平方。
\end{exercise}
\begin{exercise}{11}{基础·两点等距轨迹}
求与 $A(0,0,0)$、$B(2,0,0)$ 等距的空间动点 $M(x,y,z)$ 的轨迹方程。
\end{exercise}
\begin{exercise}{12}{常规·中位线向量证明}
三角形 $ABC$ 中 $M,N$ 分别为 $AB,AC$ 的中点。用 $\overrightarrow{MN}=\overrightarrow{AN}-\overrightarrow{AM}$ 证明 $\overrightarrow{MN}=\frac12\overrightarrow{BC}$。
\end{exercise}
\begin{exercise}{13}{基础·中点坐标变式}
求 $A(1,-2,3)$、$B(3,4,-1)$ 的中点坐标，并求原点到该中点的距离。
\end{exercise}
\begin{exercise}{14}{常规·线段交点小综合}
已知 $A(0,0,0)$、$B(2,2,0)$、$C(0,2,0)$、$D(2,0,0)$。写出直线 $AB,CD$ 的参数方程并求交点。
\end{exercise}
\subsection{内积、向量积与混合积}
\begin{exercise}{15}{基础·三种向量乘积}
设 $\vec a=(1,0,0)$、$\vec b=(0,2,0)$、$\vec c=(0,0,3)$。计算 $\vec a\cdot\vec b$、$\vec a\times\vec b$ 及 $(\vec a\times\vec b)\cdot\vec c$。
\end{exercise}
\begin{exercise}{16}{基础·混合积与体积}
设 $\vec a=(1,0,1)$、$\vec b=(0,1,1)$、$\vec c=(0,0,2)$。求三向量张成的平行六面体体积。
\end{exercise}
\begin{exercise}{17}{基础·向量积运算}
化简 $(\vec a+\vec b)\times(\vec a-\vec b)$，结果用 $\vec a\times\vec b$ 表示。
\end{exercise}
\begin{exercise}{18}{基础·内积与垂直}
设 $\vec a=(1,2,-1)$、$\vec b=(2,0,2)$，判断两向量是否垂直。
\end{exercise}
\begin{exercise}{19}{基础·向量夹角}
求 $\vec a=(1,1,0)$ 与 $\vec b=(1,0,1)$ 的夹角。
\end{exercise}
\begin{exercise}{20}{基础·三角形面积}
用向量积求 $A(0,0,0)$、$B(2,0,0)$、$C(0,3,0)$ 所成三角形的面积。
\end{exercise}
\begin{exercise}{21}{常规·面积与高}
设 $A(1,0,0)$、$B(0,1,0)$、$C(0,0,1)$。求三角形 $ABC$ 的面积及从 $A$ 到直线 $BC$ 的距离。
\end{exercise}
\begin{exercise}{22}{基础·共线与内积}
设 $\vec m=(2,-3,1)$。求与 $\vec m$ 共线且满足 $\vec r\cdot\vec m=14$ 的 $\vec r$。
\end{exercise}
\begin{exercise}{23}{基础·正交参数}
设 $\vec a=(1,2,0)$、$\vec b=(k,-1,3)$。求使 $\vec a\perp\vec b$ 的 $k$。
\end{exercise}
\begin{exercise}{24}{基础·内积与向量积关系}
设 $|\vec a|=3$、$|\vec b|=4$，两向量夹角为 $\pi/6$。求 $\vec a\cdot\vec b$ 与 $|\vec a\times\vec b|$。
\end{exercise}
\begin{exercise}{25}{常规·向量和的模}
设 $|\vec a|=2$、$|\vec b|=3$，夹角为 $2\pi/3$。求 $|\vec a+\vec b|$ 与 $|\vec a-\vec b|$。
\end{exercise}
\begin{exercise}{26}{基础·共面判定}
用混合积判断 $\vec a=(1,0,1)$、$\vec b=(0,1,1)$、$\vec c=(1,1,2)$ 是否共面。
\end{exercise}
\begin{exercise}{27}{基础·乘积恒等式核验}
设 $\vec a=(1,1,0)$、$\vec b=(0,1,1)$。计算并核验 $|\vec a\times\vec b|^2+(\vec a\cdot\vec b)^2=|\vec a|^2|\vec b|^2$。
\end{exercise}
\begin{exercise}{28}{常规·余弦定理}
利用 $\overrightarrow{BC}=\overrightarrow{AC}-\overrightarrow{AB}$ 及向量模的平方，证明三角形的余弦定理 $a^2=b^2+c^2-2bc\cos A$。
\end{exercise}
\begin{exercise}{29}{基础·两向量张成面积}
设 $\vec a=(1,0,1)$、$\vec b=(0,2,0)$。求以它们为邻边的平行四边形面积。
\end{exercise}
\begin{exercise}{30}{基础·四面体体积变式}
求以 $O(0,0,0)$、$A(2,0,0)$、$B(0,3,0)$、$C(0,0,4)$ 为顶点的四面体体积。
\end{exercise}
\begin{exercise}{31}{常规·混合积参数小综合}
设 $\vec a=(1,1,0)$、$\vec b=(0,1,1)$、$\vec c=(1,t,1)$。求使三向量共面的 $t$，并在此时将 $\vec c$ 表示为 $\vec a,\vec b$ 的线性组合。
\end{exercise}
\subsection{平面方程}
\begin{exercise}{32}{基础·点法式平面}
求过 $A(1,2,3)$、法向量为 $(2,1,-1)$ 的平面方程。
\end{exercise}
\begin{exercise}{33}{基础·三点确定平面}
求过 $A(1,0,0)$、$B(0,2,0)$、$C(0,0,3)$ 的平面方程。
\end{exercise}
\begin{exercise}{34}{基础·平行坐标轴的平面}
求平行于 $y$ 轴且过 $A(1,0,1)$、$B(2,-1,3)$ 的平面方程。
\end{exercise}
\begin{exercise}{35}{基础·特殊平面}
写出平面 $3y-1=0$ 的一个法向量，并指出它平行于哪个坐标面。
\end{exercise}
\begin{exercise}{36}{常规·垂直平面}
求过 $A(0,0,0)$、$B(1,1,0)$ 且垂直于 $x+y+z=1$ 的平面方程。
\end{exercise}
\begin{exercise}{37}{基础·双重垂直条件}
求过原点且同时垂直于平面 $x=1$、$y=2$ 的平面方程。
\end{exercise}
\begin{exercise}{38}{基础·垂直平分面}
求 $A(1,2,3)$、$B(3,2,1)$ 的垂直平分面。
\end{exercise}
\begin{exercise}{39}{基础·两方向确定平面}
求过原点且平行于向量 $(1,1,0)$、$(0,1,1)$ 的平面方程。
\end{exercise}
\begin{exercise}{40}{基础·截距式}
求在 $x,y,z$ 轴上的截距分别为 $2,3,6$ 的平面方程，并求原点到该平面的距离。
\end{exercise}
\begin{exercise}{41}{基础·平面夹角}
求平面 $x+y=0$ 与 $x-y=1$ 的夹角。
\end{exercise}
\begin{exercise}{42}{基础·坐标面夹角}
求平面 $x+y+\sqrt2z=1$ 与 $xOy$ 面的夹角。
\end{exercise}
\begin{exercise}{43}{基础·过点的平行平面}
求过 $P(1,0,0)$ 且平行于 $2x+y-z=3$ 的平面方程。
\end{exercise}
\begin{exercise}{44}{基础·平面位置关系}
判断 $x+2y-z=1$ 与 $2x+4y-2z=3$ 的位置关系；若平行，求两平面距离。
\end{exercise}
\begin{exercise}{45}{基础·平面交线}
求平面 $x+y=2$ 与 $x=1$ 的交线参数方程。
\end{exercise}
\begin{exercise}{46}{基础·平面法向量}
求平面 $2x+y-2z=5$ 与三个坐标面的夹角余弦。
\end{exercise}
\begin{exercise}{47}{常规·点面距离与垂足}
求 $P(2,0,0)$ 到平面 $x+y+z=1$ 的距离及垂足坐标。
\end{exercise}
\begin{exercise}{48}{基础·等距平行平面}
求与平面 $z=1$ 平行且相距 $2$ 的两个平面方程。
\end{exercise}
\begin{exercise}{49}{基础·简单切平面}
求曲面 $z=x^2+y^2$ 在 $(1,1,2)$ 处的切平面。
\end{exercise}
\begin{exercise}{50}{基础·截距与体积}
求平面 $x/2+y/4-z/3=1$ 与三个坐标面围成四面体的体积。
\end{exercise}
\begin{exercise}{51}{基础·平行平面距离}
求 $3x+6y-2z-7=0$ 与 $3x+6y-2z+14=0$ 的距离。
\end{exercise}
\begin{exercise}{52}{常规·球面与平面相切}
球心为 $O(1,1,1)$ 的球面与平面 $x+y+z=0$ 相切。求球面半径与切点。
\end{exercise}
\begin{exercise}{53}{常规·过直线的平面}
求包含直线 $(x,y,z)=(1,0,0)+t(1,1,0)$ 且平行于向量 $(0,0,1)$ 的平面方程。
\end{exercise}
\begin{exercise}{54}{常规·角平分面}
求平面 $x-y=0$ 与 $x+y=2$ 的两个角平分面，提示：令点到两平面的距离相等。
\end{exercise}
\begin{exercise}{55}{基础·平面束变式}
求过直线 $x=y=0$ 且过点 $P(1,2,3)$ 的平面方程。
\end{exercise}
\begin{exercise}{56}{常规·平面与截体小综合}
平面过 $P(1,1,1)$，且在三个正半轴上的截距相等。求平面方程及它与三个坐标面围成四面体的体积。
\end{exercise}
\subsection{空间直线}
\begin{exercise}{57}{基础·两点式直线}
求过 $A(1,2,3)$、$B(3,2,1)$ 的直线参数方程，注意方向向量的 $y$ 分量为零。
\end{exercise}
\begin{exercise}{58}{基础·平面法线}
求过 $P(1,0,2)$ 且垂直于 $x-2y+z=3$ 的直线参数方程。
\end{exercise}
\begin{exercise}{59}{基础·一般式转参数式}
将直线 $x+y=1$、$z=2$ 写成参数方程。
\end{exercise}
\begin{exercise}{60}{基础·公共平行方向}
求过原点且方向同时平行于平面 $x+y=1$ 和 $y+z=2$ 的直线参数方程。
\end{exercise}
\begin{exercise}{61}{基础·公共垂直方向}
求过 $P(1,1,1)$ 且方向同时垂直于 $(1,1,0)$ 与 $(0,1,1)$ 的直线参数方程。
\end{exercise}
\begin{exercise}{62}{基础·直线方向角}
求直线 $(x,y,z)=(1,2,0)+t(1,1,\sqrt2)$ 与 $z$ 轴的夹角。
\end{exercise}
\begin{exercise}{63}{常规·点到直线的垂足}
求 $P(1,0,0)$ 在直线 $(x,y,z)=t(1,1,0)$ 上的垂足及点线距离。
\end{exercise}
\begin{exercise}{64}{基础·点线距离}
求 $P(1,2,3)$ 到 $z$ 轴的距离及垂足坐标。
\end{exercise}
\begin{exercise}{65}{基础·两直线夹角}
求方向向量分别为 $(1,0,1)$、$(0,1,1)$ 的两条直线的夹角。
\end{exercise}
\begin{exercise}{66}{常规·一般式直线方向}
求直线 $x+y+z=0$、$x-y=0$ 的一个方向向量，并求该直线与 $z$ 轴的夹角。
\end{exercise}
\begin{exercise}{67}{基础·线面夹角}
求直线 $(x,y,z)=t(1,0,1)$ 与平面 $z=0$ 的夹角。
\end{exercise}
\begin{exercise}{68}{基础·方向正交参数}
两直线方向向量为 $(1,t,2)$ 和 $(4,2,-1)$。求使两方向垂直的 $t$。
\end{exercise}
\begin{exercise}{69}{基础·线面交点}
求直线 $(x,y,z)=(1,0,0)+t(1,1,1)$ 与平面 $x+y+z=4$ 的交点。
\end{exercise}
\begin{exercise}{70}{基础·直线包含于平面}
判断直线 $(x,y,z)=(1,0,0)+t(1,-1,0)$ 是否位于平面 $x+y=1$ 内。
\end{exercise}
\begin{exercise}{71}{基础·代入参数求交点}
求直线 $(x,y,z)=(0,1,2)+t(2,-1,1)$ 与平面 $x+y+z=6$ 的交点。
\end{exercise}
\begin{exercise}{72}{基础·线面平行判定}
判断直线 $(x,y,z)=(0,0,1)+t(1,-1,0)$ 与平面 $x+y+z=0$ 的位置关系。
\end{exercise}
\begin{exercise}{73}{常规·平面内指定方向直线}
求经过原点、位于平面 $x+y+z=0$ 内且方向垂直于 $(1,-1,0)$ 的直线方程。
\end{exercise}
\begin{exercise}{74}{基础·线面距离}
求直线 $(x,y,z)=(0,0,2)+t(1,1,0)$ 到平面 $z=0$ 的距离。
\end{exercise}
\begin{exercise}{75}{基础·平行线面距离}
求直线 $(x,y,z)=(1,0,0)+t(1,-1,0)$ 与平面 $x+y+z=3$ 的距离。
\end{exercise}
\begin{exercise}{76}{常规·异面直线距离}
求 $L_1:(x,y,z)=(t,0,0)$ 与 $L_2:(x,y,z)=(0,s,2)$ 的距离，并指出一对垂足。
\end{exercise}
\begin{exercise}{77}{常规·点的平面投影}
求 $P(1,2,3)$ 在平面 $x+y+z=0$ 上的投影点及关于该平面的对称点。
\end{exercise}
\begin{exercise}{78}{常规·点的直线投影}
求 $P(1,0,1)$ 在直线 $(x,y,z)=t(1,1,1)$ 上的投影点及点线距离。
\end{exercise}
\begin{exercise}{79}{基础·平行与包含}
判断直线 $(x,y,z)=(1,0,0)+t(1,-1,0)$ 与平面 $x+y=1$、$x+y=2$ 的位置关系。
\end{exercise}
\begin{exercise}{80}{常规·简单公垂线}
设 $L_1:(x,y,z)=(t,0,0)$、$L_2:(x,y,z)=(0,1,s)$。求两直线的公垂线与距离。
\end{exercise}
\begin{exercise}{81}{基础·直线夹角变式}
求直线 $(x,y,z)=t(1,1,0)$ 与 $(x,y,z)=(0,0,1)+s(1,-1,0)$ 的夹角。
\end{exercise}
\begin{exercise}{82}{常规·位置距离小综合}
设 $L_1:(x,y,z)=(t,0,0)$、$L_2:(x,y,z)=(0,1,2)+s(0,1,1)$。判断两直线是否异面，并求它们的距离。
\end{exercise}
\subsection{曲面及其方程}
\begin{exercise}{83}{基础·配方识别球面}
求球面 $x^2+y^2+z^2-6x+8y+2z+10=0$ 的球心与半径。
\end{exercise}
\begin{exercise}{84}{基础·直径端点求球面}
球面一条直径的端点为 $A(2,-3,5)$、$B(4,1,-3)$。求球面方程。
\end{exercise}
\begin{exercise}{85}{常规·四点求球面}
求过 $O(0,0,0)$、$A(2,0,0)$、$B(0,2,0)$、$C(0,0,2)$ 的球面方程。
\end{exercise}
\begin{exercise}{86}{基础·标准圆柱面}
求以 $z$ 轴为轴线、半径为 $2$ 的圆柱面方程，并写出它与 $z=1$ 的截线方程。
\end{exercise}
\begin{exercise}{87}{基础·二次曲面识别}
识别 $x^2+y^2-z^2=1$、$x^2-y^2-z^2=1$ 的曲面类型，并指出它们的轴向。
\end{exercise}
\begin{exercise}{88}{基础·旋转球面}
将 $xOz$ 面内的圆 $x^2+z^2=9$ 绕 $z$ 轴旋转一周，求所得曲面方程。
\end{exercise}
\begin{exercise}{89}{基础·旋转双曲面}
将 $xOy$ 面内的双曲线 $x^2-y^2=1$ 绕 $y$ 轴旋转，求所得曲面方程并识别类型。
\end{exercise}
\begin{exercise}{90}{基础·缺失变量与柱面}
说明空间中的 $y^2=x$、$x^2+y^2=1$ 各是什么柱面，并指出母线方向。
\end{exercise}
\begin{exercise}{91}{基础·标准二次曲面}
将 $x^2+2y^2+3z^2=6$ 化为标准形式，指出曲面类型及三个半轴长。
\end{exercise}
\begin{exercise}{92}{常规·球面相切}
求半径为 $2$ 且与平面 $z=0$ 相切于原点的两个球面方程。
\end{exercise}
\begin{exercise}{93}{常规·旋转椭球}
将 $xOy$ 面内的椭圆 $x^2/4+y^2/3=1$ 绕 $x$ 轴旋转。求椭球面方程，并求它与平面 $x=1$ 的截线半径。
\end{exercise}
\begin{exercise}{94}{基础·抛物面变式}
识别 $z=x^2+2y^2$ 的曲面类型，指出顶点，并写出 $z=2$ 的截线方程。
\end{exercise}
\begin{exercise}{95}{常规·交叉项小综合}
利用 $u=(x+z)/\sqrt2$、$v=(x-z)/\sqrt2$，将 $x^2+y^2+z^2+xz=1$ 化为标准形式，识别曲面类型。
\end{exercise}

\clearpage
\section{第9章 多元函数微分法及其应用}
\subsection{多元函数的基本概念}
\begin{exercise}{96}{基础·二元函数定义域}
求 $z=\ln(x-y)+\sqrt y$ 的定义域，并分别写出对数与根式的约束。
\end{exercise}
\begin{exercise}{97}{基础·根式分母定义域}
求 $z=\frac{\sqrt x}{\sqrt{1-x^2-y^2}}$ 的定义域，注意分母必须非零。
\end{exercise}
\begin{exercise}{98}{基础·三元函数定义域}
求 $u=\sqrt{4-x^2-y^2-z^2}$ 的定义域，并说明其几何意义。
\end{exercise}
\begin{exercise}{99}{基础·函数复合}
设 $f(u,v)=u^2+uv$，求 $f(x+y,x-y)$ 并化简。
\end{exercise}
\begin{exercise}{100}{常规·变量代换恢复函数}
已知对所有实数 $x,y$ 有 $f(x+y,x-y)=x^2+y^2$。令 $u=x+y$、$v=x-y$，求 $f(u,v)$。
\end{exercise}
\begin{exercise}{101}{基础·对数复合}
设 $f(u,v)=u+e^v$。求 $f(x-y,\ln x)$ 的表达式及定义域。
\end{exercise}
\begin{exercise}{102}{基础·有界因子极限}
用夹逼定理求 $\lim_{(x,y)\to(0,0)}(x^2+y^2)\sin\frac1{\sqrt{x^2+y^2}}$。
\end{exercise}
\begin{exercise}{103}{常规·路径判别}
沿 $y=0$ 与 $y=x$ 两条路径，说明 $\lim_{(x,y)\to(0,0)}\frac{xy}{x^2+y^2}$ 不存在。
\end{exercise}
\begin{exercise}{104}{基础·等价无穷小}
求 $\lim_{(x,y)\to(0,0)}\frac{\sin(x^2+y^2)}{x^2+y^2}$。
\end{exercise}
\begin{exercise}{105}{基础·径向极限}
求 $\lim_{(x,y)\to(0,0)}\frac{\ln(1+x^2+y^2)}{x^2+y^2}$。
\end{exercise}
\begin{exercise}{106}{基础·无穷远二重极限}
求 $x,y$ 同时趋于 $+\infty$ 时 $\frac1{x+y}$ 的极限。
\end{exercise}
\begin{exercise}{107}{基础·指数衰减}
求 $x,y$ 同时趋于 $+\infty$ 时 $e^{-x}+e^{-y}$ 的极限。
\end{exercise}
\begin{exercise}{108}{常规·幂指型极限}
求 $\lim_{(x,y)\to(0,0)}(1+x^2+y^2)^{1/(x^2+y^2)}$。
\end{exercise}
\begin{exercise}{109}{基础·连续函数代入}
求 $f(x,y)=\frac{\sqrt{4-x^2-y^2}}{1+x^2+y^2}$ 的定义域及 $\lim_{(x,y)\to(1,0)}f(x,y)$。
\end{exercise}
\begin{exercise}{110}{常规·连续延拓}
设非原点处 $f(x,y)=\frac{xy}{\sqrt{x^2+y^2}}$，$f(0,0)=a$。求使 $f$ 在原点连续的 $a$。提示：$|xy|\le(x^2+y^2)/2$。
\end{exercise}
\begin{exercise}{111}{基础·对数径向连续性}
设非原点处 $f(x,y)=(x^2+y^2)\ln(x^2+y^2)$，$f(0,0)=0$。判断 $f$ 在原点是否连续。
\end{exercise}
\begin{exercise}{112}{基础·极限估计变式}
求 $\lim_{(x,y)\to(0,0)}\frac{x^2y}{x^2+y^2}$，提示：绝对值不超过 $|y|$。
\end{exercise}
\begin{exercise}{113}{常规·连续性小综合}
设非原点处 $f(x,y)=\frac{x^2y^2}{x^2+y^2}$，$f(0,0)=a$。用 $x^2y^2\le(x^2+y^2)^2/4$ 求使 $f$ 在原点连续的 $a$。
\end{exercise}
\subsection{偏导数}
\begin{exercise}{114}{基础·偏导数定义}
已知 $f(x,0)=2x^2+3x+1$。用定义求 $f_x(0,0)$。
\end{exercise}
\begin{exercise}{115}{基础·原点偏导数}
设 $f(x,y)=\sqrt{x^2+y^4}$。用定义判断原点两个偏导数是否存在。
\end{exercise}
\begin{exercise}{116}{基础·差商换元}
设 $f_x(x_0,y_0)=3$。求 $\lim_{h\to0}\frac{f(x_0+2h,y_0)-f(x_0,y_0)}h$。
\end{exercise}
\begin{exercise}{117}{常规·不对称差商}
设 $f_x(x_0,y_0)=2$。求 $\lim_{h\to0}\frac{f(x_0+2h,y_0)-f(x_0-3h,y_0)}h$。
\end{exercise}
\begin{exercise}{118}{基础·绝对值函数偏导}
设 $f(x,y)=|x|y$。用定义求 $f_x(0,0)$、$f_y(0,0)$。
\end{exercise}
\begin{exercise}{119}{基础·根式偏导数}
设 $f(x,y)=\sqrt{|x|}+y^2$。用定义判断原点两个偏导数是否存在。
\end{exercise}
\begin{exercise}{120}{基础·多项式偏导数}
设 $f(x,y)=x^3+3xy+4y-1$。求 $f_x(1,2)$、$f_y(1,2)$。
\end{exercise}
\begin{exercise}{121}{基础·反三角函数偏导}
设 $f(x,y)=x+2y+(y-1)\arcsin x$，$|x|<1$。求 $f_x(0,1)$、$f_y(0,1)$。
\end{exercise}
\begin{exercise}{122}{基础·立方根偏导定义}
设 $f(x,y)=\sqrt[3]{x^3-y^3}$，立方根取实值。用定义求原点的两个偏导数。
\end{exercise}
\begin{exercise}{123}{基础·指数三角偏导}
求 $z=e^x\sin(x^2+y)$ 的两个一阶偏导数。
\end{exercise}
\begin{exercise}{124}{基础·三元幂函数偏导}
设 $u=(x/y)^z$，其中 $x>0,y>0,z\in\mathbb R$。求三个一阶偏导数。
\end{exercise}
\begin{exercise}{125}{常规·变底变指数}
在 $x+y>0$ 上设 $z=(x+y)^x$。求 $z_x,z_y$ 及它们在 $(1,1)$ 处的值。
\end{exercise}
\begin{exercise}{126}{常规·对数求导}
设 $z=(y/x)^{x/y}$，$x,y>0$。求 $z_x(1,1)$、$z_y(1,1)$。
\end{exercise}
\begin{exercise}{127}{基础·分段函数偏导}
设非原点处 $f(x,y)=(x^2+y^2)\ln(x^2+y^2)$，原点处取零。用定义求 $f_x(0,0)$、$f_y(0,0)$。
\end{exercise}
\begin{exercise}{128}{常规·由偏导恢复函数}
设 $x>0$，$f_x=x+y^2/x$，$f(1,y)=\sin y$。求 $f(x,y)$。
\end{exercise}
\begin{exercise}{129}{基础·截线斜率}
设曲面 $z=\sqrt{1+x^2+y^2}$。求固定 $x=1$ 的截线在 $(1,1,\sqrt3)$ 处关于 $y$ 的导数。
\end{exercise}
\begin{exercise}{130}{基础·截线切线}
求曲线 $z=(x^2+y^2)/4$、$y=4$ 在 $(2,4,5)$ 处的切线参数方程。
\end{exercise}
\begin{exercise}{131}{常规·偏导与单调性}
设 $f(x,y)=e^x-y$。利用两个偏导数的符号，比较 $f(0,2)$ 与 $f(1,1)$ 的大小。
\end{exercise}
\begin{exercise}{132}{基础·连续与偏导变式}
设 $f(x,y)=\sqrt{x^2+y^2}$。判断它在原点的连续性及两个偏导数是否存在。
\end{exercise}
\begin{exercise}{133}{常规·偏导与连续性辨析}
设原点处 $f=0$，非原点处 $f(x,y)=\frac{2xy}{x^2+y^2}$。求原点两个偏导数，并沿 $y=x$ 判断原点连续性。
\end{exercise}
\subsubsection{高阶偏导数}
\begin{exercise}{134}{基础·二阶偏导}
设 $z=x^3y+xy^2$。求 $z_{xx},z_{xy},z_{yy}$。
\end{exercise}
\begin{exercise}{135}{常规·积分上限二阶偏导}
设 $g(x,y)=\int_0^{xy}\sin t\dd t$。求 $g_x$ 和 $g_{xx}(0,2)$。
\end{exercise}
\begin{exercise}{136}{基础·混合高阶偏导}
设 $z=e^{x+y}$。求 $z_{xxy}$ 及其在原点的值。
\end{exercise}
\begin{exercise}{137}{常规·径向二阶偏导}
在 $r=\sqrt{x^2+y^2}>0$ 上设 $u=r$。求 $u_{xx},u_{yy}$，并验证 $u_{xx}+u_{yy}=1/r$。
\end{exercise}
\begin{exercise}{138}{基础·混合偏导交换}
设 $f(x,y)=x^2y^3$。分别计算 $f_{xy},f_{yx}$，核验二者相等。
\end{exercise}
\begin{exercise}{139}{基础·二阶偏导变式}
设 $f(x,y)=\sin(x+y)$。求原点的三个二阶偏导数 $f_{xx},f_{xy},f_{yy}$。
\end{exercise}
\begin{exercise}{140}{常规·二阶偏导小综合}
设 $f(x,y)=e^{xy}$。求 $f_{xx},f_{xy},f_{yy}$ 及它们在原点的值。
\end{exercise}
\subsection{全微分}
\begin{exercise}{141}{基础·全增量与全微分}
设 $f(x,y)=x^2+y^2$。求在 $(1,2)$ 处当 $\Delta x=0.1$、$\Delta y=-0.1$ 时的全增量和全微分。
\end{exercise}
\begin{exercise}{142}{基础·全微分计算}
求 $z=(x^2+y^2)e^x$ 的全微分。
\end{exercise}
\begin{exercise}{143}{基础·三元函数全微分}
设 $u=\ln(x+y+z)$，$x+y+z>0$。求 $\dd u$。
\end{exercise}
\begin{exercise}{144}{常规·隐式全微分}
设 $z-x-y+xe^z=0$ 在原点附近确定 $z=z(x,y)$。求 $\dd z|_{(0,0)}$。
\end{exercise}
\begin{exercise}{145}{基础·原点可微性}
设 $f(x,y)=\frac{2x+3y}{1+x^2+y^2}$。求原点的全微分。
\end{exercise}
\begin{exercise}{146}{常规·渐近式与全微分}
设 $f$ 在原点连续，且 $f(x,y)=1+x+2y+o(\sqrt{x^2+y^2})$。求 $f(0,0)$ 及原点全微分。
\end{exercise}
\begin{exercise}{147}{常规·全微分恢复数据}
设 $u\in C^2$ 且 $u_x=3x^2y$，$u(0,y)=y^2$。求 $u(x,y)$ 与 $u_y$。
\end{exercise}
\begin{exercise}{148}{常规·由全微分恢复函数}
已知 $\dd f=2xy\dd x+(x^2+2y)\dd y$，$f(0,0)=2$。求 $f(x,y)$。
\end{exercise}
\begin{exercise}{149}{基础·可微定义}
设 $f(x,y)=3x-2y+x^2+y^2$。用可微定义求原点全微分，提示：令 $r=\sqrt{x^2+y^2}$，考察余项与 $r$ 的比值。
\end{exercise}
\begin{exercise}{150}{常规·不连续与不可微}
设原点处 $f=0$，非原点处 $f(x,y)=\frac{3x^2y}{x^4+y^2}$。沿 $y=x^2$ 判断连续性，从而判断原点可微性。
\end{exercise}
\begin{exercise}{151}{基础·根式函数可微性}
判断 $f(x,y)=\sqrt{x^2+y^2}$ 在原点是否可微，提示：先检查偏导数。
\end{exercise}
\begin{exercise}{152}{常规·振荡函数可微性}
设原点处 $f=0$，非原点处 $f(x,y)=xy\sin\frac1{\sqrt{x^2+y^2}}$。用 $|xy|\le(x^2+y^2)/2$ 证明它在原点可微并求全微分。
\end{exercise}
\begin{exercise}{153}{基础·可微充分条件}
设 $f(x,y)=\ln(1+x^2+y^2)$。说明它在原点附近的一阶偏导数连续，并求原点全微分。
\end{exercise}
\begin{exercise}{154}{常规·连续但不可微}
设原点处 $f=0$，非原点处 $f(x,y)=\frac{x^2y}{x^2+y^2}$。求原点两个偏导数，并沿 $y=x$ 检查可微定义中的归一化余项。
\end{exercise}
\begin{exercise}{155}{常规·可微性辨析}
设原点处 $f=0$，非原点处 $f(x,y)=\frac{xy}{\sqrt{x^2+y^2}}$。求原点两个偏导数，并沿 $y=x$ 判断可微性。
\end{exercise}
\begin{exercise}{156}{常规·径向余项估计}
设原点处 $f=0$，非原点处 $f(x,y)=(x^2+y^2)\sin\frac1{x^2+y^2}$。用可微定义求原点全微分。
\end{exercise}
\begin{exercise}{157}{基础·可微性变式}
设 $f(x,y)=1+2x-y+(x^2+y^2)^{3/2}$。用可微定义求原点全微分。
\end{exercise}
\begin{exercise}{158}{常规·绝对值小综合}
设 $f(x,y)=|xy|$。用定义求原点两个偏导数，并利用 $|xy|\le(x^2+y^2)/2$ 判断可微性。
\end{exercise}
\subsection{多元复合函数的求导法则}
\begin{exercise}{159}{基础·一元链式法则}
设 $z=u^2+v^2$，$u=e^x$、$v=x^2$。用链式法则求 $\dd z/\dd x$。
\end{exercise}
\begin{exercise}{160}{基础·全导数}
设 $f(x,y)=xy$，$z(x)=f(x,\ln x)$，$x>0$。求 $z'(x)$ 及 $z'(1)$。
\end{exercise}
\begin{exercise}{161}{基础·抽象函数链式法则}
设 $z=f(x^2,\sin x)$，$f\in C^1$。求 $z'(x)$，写明 $f_1,f_2$ 的取值点。
\end{exercise}
\begin{exercise}{162}{常规·二阶链式法则}
设 $z=f(x,x^2)$，$f\in C^2$。求 $z''(0)$，结果用 $f$ 在原点的一、二阶偏导数表示。
\end{exercise}
\begin{exercise}{163}{常规·驻点内函数}
设 $f\in C^2$，$f_1(1,1)=3$、$f_2(1,1)=4$，$z=f(\cos x,1+x^2)$。求 $z'(0),z''(0)$。
\end{exercise}
\begin{exercise}{164}{基础·两个中间变量}
设 $z=e^u\sin v$，$u=x-y$、$v=x+y$。求 $z_x,z_y$。
\end{exercise}
\begin{exercise}{165}{基础·径向复合求导}
设 $z=\ln(1+\sqrt{x^2+y^2})$。求原点外的 $z_x,z_y$。
\end{exercise}
\begin{exercise}{166}{基础·显含变量链式法则}
设 $z=x^2+y^2+u^2$，$u=xy$。求 $z_x,z_y$。
\end{exercise}
\begin{exercise}{167}{基础·抽象复合偏导}
设 $z=f(x+y,xy)$，$f\in C^1$。求 $z_x,z_y$ 及它们在原点的表达式。
\end{exercise}
\begin{exercise}{168}{常规·线性变换二阶偏导}
设 $z=f(x+y,x-y)$，$f\in C^2$。求 $z_{xx},z_{yy}$，明确外函数偏导数的取值点。
\end{exercise}
\begin{exercise}{169}{常规·混合偏导特殊点}
设 $z=f(e^{xy},x^2+y^2)$，$f\in C^2$。求 $z_{xy}(0,0)$，结果用 $f$ 在 $(1,0)$ 的偏导数表示。
\end{exercise}
\begin{exercise}{170}{常规·截面条件恢复函数}
设 $z=e^{x-y}+f(x-2y)$，且 $z(x,0)=x^2$。求 $f$ 及 $z_y$。
\end{exercise}
\begin{exercise}{171}{基础·特殊点全微分}
设 $z=yf(x^2-y^2)$，$f$ 可导，$f(0)=1$、$f'(0)=1/2$。求 $\dd z|_{(1,1)}$。
\end{exercise}
\begin{exercise}{172}{基础·变限积分偏导}
设 $z=\int_0^{x^2+y}e^{-t^2}\dd t$。求 $z_x,z_y$。
\end{exercise}
\begin{exercise}{173}{常规·齐次恒等式}
设 $z=xf(y/x)$，$x\ne0$，$f$ 可导。直接求 $z_x,z_y$，验证 $xz_x+yz_y=z$。
\end{exercise}
\begin{exercise}{174}{基础·复合函数偏导}
设 $z=xf(x/y)$，$y>0$，$f$ 可导。求 $z_x,z_y$。
\end{exercise}
\begin{exercise}{175}{常规·混合偏导计算}
设 $z=f(2x-y)+x^2y$，$f$ 二阶可导。求 $z_{xx},z_{xy}$。
\end{exercise}
\begin{exercise}{176}{基础·三元外函数复合}
设 $u=f(x,y,xe^y)$，$f\in C^1(\mathbb R^3)$。求 $u_x,u_y$。
\end{exercise}
\begin{exercise}{177}{常规·多层链式法则}
设 $u=f(x,y,z)$，其中 $y=x+t$、$t=xz$，$x,z$ 为独立变量，$f\in C^1$。求复合函数的 $u_x,u_z$。
\end{exercise}
\begin{exercise}{178}{常规·线性代换混合偏导}
设 $z=g(x+y,x-y)$，$g\in C^2$。求 $z_{xy}$，结果用 $g$ 的二阶偏导数表示。
\end{exercise}
\begin{exercise}{179}{基础·齐次公式核验}
设 $f(x,y)=x^2+xy+y^2$。验证 $xf_x+yf_y=2f$。
\end{exercise}
\begin{exercise}{180}{基础·全微分形式不变性}
设 $z=\arctan(u+v)$，$u=2x-y^2$、$v=x^2y$。求 $(1,0)$ 处的全微分。
\end{exercise}
\begin{exercise}{181}{基础·复合全微分}
设 $z=f(2x+3y,e^{xy})$，$f\in C^1$。求原点处的全微分，结果用 $f_1(0,1),f_2(0,1)$ 表示。
\end{exercise}
\begin{exercise}{182}{基础·链式法则变式}
设 $F(u,v)=u^2+v^2$，$u=t$、$v=2t$。用链式法则求 $\dd F/\dd t$。
\end{exercise}
\begin{exercise}{183}{常规·极坐标求导小综合}
设 $x=r\cos\theta$、$y=r\sin\theta$、$z=x^2+2y^2$。求 $z_r,z_\theta$，并计算它们在 $(r,\theta)=(1,\pi/4)$ 处的值。
\end{exercise}
\subsection{隐函数的求导法则}
\begin{exercise}{184}{基础·隐函数一阶导数}
设 $e^x+y^2=2$ 在 $(0,1)$ 附近确定 $y=y(x)$。求 $y'$ 及 $y'(0)$。
\end{exercise}
\begin{exercise}{185}{基础·隐函数初值}
设 $2^{xy}=x+y$ 在 $x=0$ 附近确定 $y=y(x)$，所需隐函数条件已满足。求 $y(0),y'(0)$。
\end{exercise}
\begin{exercise}{186}{基础·二阶隐式求导}
设 $x^2+y^2=1$ 在 $(0,1)$ 附近确定正值分支 $y=y(x)$。求 $y'(0),y''(0)$。
\end{exercise}
\begin{exercise}{187}{常规·含未知函数的隐式求导}
设 $f\in C^1$，$xf(x+y)=\ln y$ 在 $(0,1)$ 附近确定 $y=y(x)$。求 $y'(0)$，结果可含 $f(1)$。
\end{exercise}
\begin{exercise}{188}{常规·隐函数复合求导}
设 $y+xe^y=1$ 在 $x=0$ 附近确定 $y=y(x)$，$z=y^2+x$。求 $z'(0)$。
\end{exercise}
\begin{exercise}{189}{基础·隐函数两个偏导}
设 $x^2+y^2+z^2=1$ 在 $(0,0,1)$ 附近确定 $z=z(x,y)$。求 $z_x,z_y$ 及原点处的值。
\end{exercise}
\begin{exercise}{190}{基础·指数隐函数}
设 $z+e^z=x+y+1$ 在原点附近确定 $z=z(x,y)$。求 $z_x(0,0),z_y(0,0)$。
\end{exercise}
\begin{exercise}{191}{基础·分式隐函数偏导}
设 $z+xz+yz=1$ 在 $1+x+y\ne0$ 的区域上确定 $z=z(x,y)$。求 $z_x,z_y$。
\end{exercise}
\begin{exercise}{192}{常规·共同平移结构}
设 $(x-y)^3+z+\tan z=0$ 在 $\tan z$ 有定义的点附近确定 $z=z(x,y)$。求 $z_x+z_y$。
\end{exercise}
\begin{exercise}{193}{基础·隐式全微分}
设 $e^z-xy^2+\sin(y+z)=0$ 确定 $z=z(x,y)$，已知 $e^z+\cos(y+z)\ne0$。求 $\dd z$。
\end{exercise}
\begin{exercise}{194}{常规·隐函数全微分}
设 $xyz+\sqrt{x^2+y^2+z^2}=\sqrt2$ 在 $(1,0,1)$ 附近确定 $z=z(x,y)$。求 $\dd z|_{(1,0)}$，所需隐函数条件已满足。
\end{exercise}
\begin{exercise}{195}{常规·隐函数二阶偏导}
设 $z^3-3xyz=1$ 在 $(0,0,1)$ 附近确定 $z=z(x,y)$。求 $z_{xy}(0,0)$。
\end{exercise}
\begin{exercise}{196}{基础·隐函数特殊点偏导}
设 $z+e^x-y\ln(1+z^2)=0$ 在 $(x,y)=(0,0)$ 附近确定 $z=z(x,y)$。求对应 $z$ 值及 $z_x(0,0),z_y(0,0)$。
\end{exercise}
\begin{exercise}{197}{常规·隐式中间变量}
设 $z+x+y=0$，$F(x,y)=yz^2e^x$。用链式法则求 $F_x(0,1)$。
\end{exercise}
\begin{exercise}{198}{基础·隐式全微分变式}
设 $z^2+x^2+y^2=4$ 在 $(0,0,2)$ 附近确定 $z=z(x,y)$。求原点处的全微分。
\end{exercise}
\begin{exercise}{199}{常规·隐函数微分小综合}
设 $z+xyz=1$ 在 $(0,0,1)$ 附近确定 $z=z(x,y)$。求 $\dd z|_{(0,0)}$ 及 $z_{xy}(0,0)$。
\end{exercise}
\subsubsection{方程组确定的隐函数}
\begin{exercise}{200}{基础·隐式方程组求导}
设 $x+y+z=1$、$y-z=x^2$ 确定 $y=y(x),z=z(x)$。对两式求导，求 $y'(0),z'(0)$。
\end{exercise}
\begin{exercise}{201}{常规·非线性方程组偏导}
设 $u^2+v=x$、$u+v^2=y$ 在 $(x,y,u,v)=(2,2,1,1)$ 附近确定 $u(x,y),v(x,y)$，雅可比条件已满足。求 $u_x(2,2),v_x(2,2)$。
\end{exercise}
\begin{exercise}{202}{常规·初值方程组}
设 $f\in C^1$、$f(-1)=1$，$z=xf(x-y)$、$x+y+z=1$ 在 $x=0$ 附近确定 $y(x),z(x)$。求 $y(0),z(0)$ 及 $y'(0),z'(0)$。
\end{exercise}
\begin{exercise}{203}{基础·隐式复合全导数}
设 $y+x^2=1$、$z+xz=1$，$u=x+y+z$。求 $u'(0)$。
\end{exercise}
\begin{exercise}{204}{基础·隐函数求导变式}
设 $z+z^3=x+y$ 在原点附近确定 $z=z(x,y)$。求 $z_x(0,0),z_y(0,0)$。
\end{exercise}
\begin{exercise}{205}{常规·隐式二阶微分小综合}
设 $z+\sin z=x+y$ 在原点附近确定 $z=z(x,y)$。求 $z_x(0,0)$ 与 $z_{xx}(0,0)$。
\end{exercise}
\subsection{多元函数微分法在几何上的应用}
\begin{exercise}{206}{基础·参数曲线切线}
设 $x=t+1$、$y=t^2-1$、$z=t^3$。求曲线在 $t=1$ 处的切线与法平面。
\end{exercise}
\begin{exercise}{207}{基础·交曲线切线}
求曲线 $y=x$、$z=x^2-2y$ 在原点的切线，提示：以 $x$ 为参数。
\end{exercise}
\begin{exercise}{208}{常规·平行切线反求点}
对曲线 $y=x^2$、$z=x^3$，求与平面 $x+y=0$ 平行的切线的切点及切线方程。
\end{exercise}
\begin{exercise}{209}{常规·交面曲线切线}
求曲线 $x^2+y^2+z^2=6$、$x+y+z=0$ 在 $(1,-2,1)$ 处的切线，提示：取两个曲面法向量的向量积。
\end{exercise}
\begin{exercise}{210}{基础·隐式曲面切平面}
求曲面 $x^2+\cos(xy)+yz+x=0$ 在 $(0,1,-1)$ 处的切平面与法线。
\end{exercise}
\begin{exercise}{211}{基础·椭球切平面}
求椭球面 $x^2+2y^2+3z^2=9$ 在 $(2,1,1)$ 处的切平面与法线。
\end{exercise}
\begin{exercise}{212}{基础·隐式曲面法线}
求曲面 $e^z-z+xy=3$ 在 $(2,1,0)$ 处的切平面与法线。
\end{exercise}
\begin{exercise}{213}{基础·以不同变量显式表示的曲面}
求曲面 $y=x^2+3z^2-15$ 在 $(1,-2,2)$ 处的切平面。
\end{exercise}
\begin{exercise}{214}{常规·平行切平面}
求抛物面 $z=4-x^2-y^2$ 上切平面平行于 $2x+2y+z=1$ 的点及对应切平面。
\end{exercise}
\begin{exercise}{215}{常规·切平面与点的判定}
求曲面 $z=x^2+y^2$ 在 $(1,-2,5)$ 处的切平面，并判断原点是否在此切平面内。
\end{exercise}
\begin{exercise}{216}{基础·球面切平面变式}
求球面 $x^2+y^2+z^2=9$ 在 $(1,2,2)$ 处的切平面。
\end{exercise}
\begin{exercise}{217}{常规·指定法向小综合}
求单位球面上切平面平行于 $x+2y+z=0$ 的两个点及对应切平面。
\end{exercise}
\subsection{方向导数和梯度}
\begin{exercise}{218}{基础·方向单位化}
求 $f(x,y)=3x^2y-y^2$ 在 $(1,1)$ 处沿向量 $(1,\sqrt3)$ 的方向导数。
\end{exercise}
\begin{exercise}{219}{常规·曲线切向方向导数}
曲线 $y=x^2-5$ 的正向为 $x$ 增大方向。求 $f(x,y)=x^2-xy+y^2$ 在 $(2,-1)$ 沿该曲线切线正向的方向导数。
\end{exercise}
\begin{exercise}{220}{基础·空间方向导数}
设 $f(x,y,z)=x^2+y^2+z^2$。求 $(1,0,1)$ 处沿 $(1,1,1)$ 的方向导数。
\end{exercise}
\begin{exercise}{221}{基础·方向导数定义}
设 $f(x,y)=x^2+y^2$。用定义 $D_{\vec e}f(0)=\lim_{t\to0^+}[f(t\vec e)-f(0)]/t$，求原点沿单位向量 $\vec e=(3/5,4/5)$ 的方向导数。
\end{exercise}
\begin{exercise}{222}{常规·法向方向导数}
设 $P(1,1,1)$，求 $u=x+y+z$ 在 $P$ 处沿球面 $x^2+y^2+z^2=3$ 外单位法向的方向导数。
\end{exercise}
\begin{exercise}{223}{基础·零方向导数}
设 $f(x,y)=x^2+y^2$。求 $(1,1)$ 处使方向导数为零的两个单位方向。
\end{exercise}
\begin{exercise}{224}{基础·最快增长方向}
设 $f(x,y,z)=x^3-xy^2-z$。求 $(1,1,0)$ 处增长最快的单位方向及最大变化率。
\end{exercise}
\begin{exercise}{225}{基础·梯度与方向导数}
求 $f(x,y,z)=x^2+y^2+xz$ 在 $(1,0,1)$ 处的梯度及沿 $(3,-1,6)$ 的方向导数。
\end{exercise}
\begin{exercise}{226}{常规·另一函数梯度方向}
设 $u=\sqrt{x^2+y^2+z^2}$、$g=x^2/2+yz$。求 $(1,1,-1)$ 处沿 $g$ 的梯度方向的 $u$ 的方向导数。
\end{exercise}
\begin{exercise}{227}{常规·隐函数梯度}
设 $x+y^2+z^3=3$ 在 $(1,1,1)$ 附近确定 $z=z(x,y)$。求 $(1,1)$ 处方向导数的最大值及取得该值的单位方向。
\end{exercise}
\begin{exercise}{228}{基础·梯度计算变式}
设 $F(x,y)=e^{x+y}$。求原点处的梯度及沿 $(1,-1)$ 的方向导数。
\end{exercise}
\begin{exercise}{229}{常规·方向导数小综合}
设温度 $T(x,y)=x^2+2y^2$。求 $(1,1)$ 处沿 $(1,1)$ 的方向导数，并求该点下降最快的单位方向。
\end{exercise}
\subsection{多元函数的极值及其应用}
\begin{exercise}{230}{基础·极值必要条件}
设 $f$ 在内点 $(x_0,y_0)$ 可微并取得局部极小值。利用两条坐标截线的一元极值必要条件，说明 $f_x(x_0,y_0)=f_y(x_0,y_0)=0$。
\end{exercise}
\begin{exercise}{231}{基础·驻点与鞍点}
求 $f(x,y)=x^2-y^2$ 的驻点，并沿坐标轴判断原点是否为极值点。
\end{exercise}
\begin{exercise}{232}{基础·平方和极值}
求 $f(x,y)=x^2+2y^2$ 的极值及极值点。
\end{exercise}
\begin{exercise}{233}{基础·四次函数极值}
求 $f(x,y)=-x^4-y^4$ 的极值及极值点。
\end{exercise}
\begin{exercise}{234}{基础·二阶极值判别}
求 $z=x^2+y^2-2x-4y$ 的驻点并判定极值。
\end{exercise}
\begin{exercise}{235}{常规·指数函数局部极值}
求 $f(x,y)=e^{-x}(x-y^2)$ 的驻点及局部极值，使用二阶判别法。
\end{exercise}
\begin{exercise}{236}{常规·指数函数极值}
求 $f(x,y)=e^{-x^2-y^2}$ 的驻点及极值。
\end{exercise}
\begin{exercise}{237}{常规·截面数据与极值}
设 $f(x,y)=2y^2+g(x-y)$，且 $f(x,0)=x^2$。求 $f$，并用配方法求其最小值。
\end{exercise}
\begin{exercise}{238}{基础·退化二阶判别}
用配方法求 $f(x,y)=(x+y)^2-2(x+y)+5$ 的最小值及取值点。
\end{exercise}
\begin{exercise}{239}{常规·隐函数极值}
设 $\ln z+x^2+y^2=0$，$z>0$。先显式写出 $z(x,y)$，再求其极值。
\end{exercise}
\begin{exercise}{240}{常规·水箱材料最省}
无盖长方体水箱的容积为 $4$ 立方米，底面边长为 $x,y>0$。写出用料面积 $S(x,y)$，求使面积最小的长、宽、高。提示：由 $xyh=4$ 消去 $h$，也可用基本不等式核验最小值。
\end{exercise}
\begin{exercise}{241}{常规·闭区域最值}
求 $f(x,y)=x^2+y^2-2x$ 在 $x^2+y^2\le4$ 上的最大值、最小值及取值点。
\end{exercise}
\begin{exercise}{242}{基础·简单条件极值}
在约束 $x+y=2$ 下求 $f(x,y)=x^2+y^2$ 的最小值及取值点。
\end{exercise}
\begin{exercise}{243}{基础·拉格朗日乘数法}
用拉格朗日乘数法求 $x^2+y^2$ 在 $x-y+1=0$ 下的最小值。
\end{exercise}
\begin{exercise}{244}{常规·球内接长方体}
长方体的中心为原点，各面平行于坐标面，八个顶点位于单位球面上。求其最大体积及三个边长。提示：第一卦限顶点满足 $x^2+y^2+z^2=1$。
\end{exercise}
\begin{exercise}{245}{常规·恢复函数再求最值}
设 $f_x=2x$、$f_y=2y$、$f(0,0)=1$。求 $f$，并求它在 $x^2+y^2\le1$ 上的最大值、最小值。
\end{exercise}
\begin{exercise}{246}{基础·退化驻点变式}
判断 $f(x,y)=x^4+y^4$ 在原点的极值类型，并求极值。
\end{exercise}
\begin{exercise}{247}{常规·圆上最值小综合}
求 $F(x,y)=x^4+y^4$ 在单位圆 $x^2+y^2=1$ 上的最大值、最小值及取值点，提示：利用 $F=1-2x^2y^2$。
\end{exercise}
\subsection{二元函数的泰勒公式}
\begin{exercise}{248}{基础·二阶泰勒多项式}
求 $f(x,y)=2x^2-xy-y^2+3x+5y-1$ 在原点的二阶泰勒多项式。
\end{exercise}
\begin{exercise}{249}{常规·非原点二阶展开}
求 $f(x,y)=x^2y$ 在 $(1,1)$ 处的二阶泰勒公式，余项用 $o(h^2+k^2)$ 表示，其中 $h=x-1,k=y-1$。
\end{exercise}
\begin{exercise}{250}{基础·麦克劳林展开}
求 $f(x,y)=e^{x+y}$ 的二阶麦克劳林公式，余项用 $o(x^2+y^2)$ 表示。
\end{exercise}
\begin{exercise}{251}{常规·泰勒展开求极限}
先写出 $\ln(1+x+y)$ 的二阶麦克劳林公式，再求 $\lim_{(x,y)\to(0,0)}\frac{\ln(1+x+y)-x-y+xy}{x^2+y^2}$。
\end{exercise}
\begin{exercise}{252}{基础·三角函数展开}
求 $f(x,y)=\cos x\cos y$ 的二阶麦克劳林公式。
\end{exercise}
\begin{exercise}{253}{基础·二阶展开变式}
求 $f(x,y)=\ln(1+x^2+y^2)$ 的二阶麦克劳林公式。
\end{exercise}
\begin{exercise}{254}{常规·展开与极值小综合}
求 $F(x,y)=e^{x+y}-1-x-y+\frac12(x-y)^2$ 的二阶麦克劳林公式，并据二次主项判定原点的极值类型。
\end{exercise}
\subsubsection{常规综合练习}
\begin{exercise}{255}{基础·偏导符号比较}
设 $f(x,y)=e^x-y$。比较 $f(0,1)$、$f(0,0)$、$f(1,0)$ 的大小，并用偏导数符号说明依据。
\end{exercise}
\begin{exercise}{256}{常规·渐近式与可微性}
设 $f$ 在原点连续，且 $\lim_{(x,y)\to(0,0)}\frac{f(x,y)-1-2x+y}{x^2+y^2}=3$。求 $f(0,0)$ 及原点全微分。
\end{exercise}
\begin{exercise}{257}{常规·有界振荡可微性}
设 $xy\ne0$ 时 $f(x,y)=(x^2+y^2)\sin(1/(xy))$，$xy=0$ 时取零。用可微定义求原点全微分。
\end{exercise}
\begin{exercise}{258}{常规·二阶渐近与极值}
设 $f$ 在原点连续，且 $\lim_{(x,y)\to(0,0)}\frac{f(x,y)-f(0,0)}{x^2+y^2}=2$。求原点两个偏导数，并判定原点的极值类型。
\end{exercise}
\begin{exercise}{259}{基础·全微分与方向导数}
已知 $f$ 在原点可微，$f_x(0,0)=1$、$f_y(0,0)=2$。求原点全微分及沿向量 $(3,4)$ 的方向导数。
\end{exercise}
\begin{exercise}{260}{常规·复合数据求导}
设 $h(x,y)=f(x-y,x+y)$，$f\in C^1$，$f_1(1,3)=2$、$f_2(1,3)=3$。求 $h_x(2,1),h_y(2,1)$。
\end{exercise}
\begin{exercise}{261}{常规·二阶偏导恢复函数}
已知 $f_{yy}=2$、$f(x,0)=1$、$f_y(x,0)=x$。求 $f(x,y)$，并判断原点是否为极值点。
\end{exercise}
\begin{exercise}{262}{常规·沿曲线求导}
设 $f\in C^2$，$f_{xx}=f_{yy}$，且 $f(x,3x)=x$、$f_x(x,3x)=0$。分别对给定关系求导，求 $f_{xx}(x,3x),f_{xy}(x,3x)$。
\end{exercise}
\begin{exercise}{263}{基础·三元外函数全微分}
设 $z=f(x,x+y,x/y)$，$f\in C^1$、$y\ne0$。求 $(0,1)$ 处的全微分，写明外函数偏导数的取值点。
\end{exercise}
\begin{exercise}{264}{常规·三元复合混合偏导}
设 $u=f(x,xy,xyz)$，$f\in C^2$。只求 $u_{xy}(0,1,1)$，结果用 $f$ 在原点的一、二阶偏导数表示。
\end{exercise}
\begin{exercise}{265}{基础·零全微分}
设 $\lim_{(x,y)\to(0,0)}\frac{f(x,y)-f(0,0)}{\sqrt{x^2+y^2}}=0$。直接根据可微定义求原点全微分。
\end{exercise}
\begin{exercise}{266}{基础·二重极限}
求 $\lim_{(x,y)\to(0,0)}\frac{\sin(x^2+y^2)}{\sqrt{1+x^2+y^2}+1}$。
\end{exercise}
\begin{exercise}{267}{常规·有理化与可微性}
设原点处 $f=0$，非原点处 $f(x,y)=\frac{\sqrt{1+x^2y^2}-1}{x^2+y^2}$。先有理化，判断原点连续性并求全微分。
\end{exercise}
\begin{exercise}{268}{常规·隐式约束全导数}
设 $u=f(x,y,z)$、$y=\sin x$、$z^2+x^2+y^2=1$，取原点附近 $z>0$ 的分支，$f\in C^1$。求 $\dd u/\dd x|_{x=0}$，结果用 $f$ 在 $(0,0,1)$ 的偏导数表示。
\end{exercise}
\begin{exercise}{269}{常规·径向函数二阶偏导}
设 $z=\ln\sqrt{x^2+y^2}$，$(x,y)\ne(0,0)$。计算 $z_{xx},z_{yy}$ 并验证 $z_{xx}+z_{yy}=0$。
\end{exercise}
\begin{exercise}{270}{常规·线性变量变换}
设 $u=x+y$、$v=x-y$、$z=g(u,v)$，$g\in C^2$。证明 $z_{xx}-z_{yy}=4g_{uv}$，提示：分别应用两次链式法则。
\end{exercise}
\begin{exercise}{271}{基础·隐式齐次关系}
设 $z=\sqrt{x^2+y^2}$，$(x,y)\ne(0,0)$。计算并验证 $xz_x+yz_y=z$。
\end{exercise}
\begin{exercise}{272}{常规·交曲线二阶导数}
设 $z=x^2+y^2$、$x+y=1$ 确定 $y=y(x),z=z(x)$。求 $y'',z''$。
\end{exercise}
\begin{exercise}{273}{常规·极值点复合二阶导数}
设 $z=f(xy,y)$，$f\in C^2$，$(1,1)$ 是 $f$ 的极值点。将 $z_{xy}(1,1)$ 用 $f$ 在 $(1,1)$ 的二阶偏导数表示。
\end{exercise}
\begin{exercise}{274}{常规·曲线数据确定切平面}
设 $f$ 可微，$f(t^2,t+1)=t^2(t-1)$，且 $f_1(1,2)=1$。求 $f_2(1,2)$ 及曲面 $z=f(x,y)$ 在对应点的切平面。
\end{exercise}
\begin{exercise}{275}{基础·指定点切平面}
求椭球面 $x^2+2y^2+3z^2=21$ 在 $(3,0,2)$ 处的切平面与法线。
\end{exercise}
\begin{exercise}{276}{常规·全微分恢复与梯度}
已知 $\dd z=(2x+y)\dd x+(x+2y)\dd y$，$z(0,0)=1$。求 $z$，并求 $(1,0)$ 处方向导数的最大值。
\end{exercise}
\begin{exercise}{277}{常规·隐函数分支极值}
方程 $z^2=1+x^2+y^2$ 确定正、负两个分支。分别求两个分支的极值及极值点。
\end{exercise}
\begin{exercise}{278}{常规·正变量约束极值}
设 $x,y,z>0$、$x+y+z=3$。求 $xyz$ 的最大值及取值点，可用拉格朗日乘数法或基本不等式。
\end{exercise}
\begin{exercise}{279}{常规·双约束最值}
在 $x+y+z=2$、$z=0$ 的约束下，求 $u=x^2+y^2+z^2$ 的最小值及取值点。
\end{exercise}
\begin{exercise}{280}{常规·方向导数与约束最值}
设空间函数 $f(x,y,z)=x^2+y^2+z^2$。在单位球面 $x^2+y^2+z^2=1$ 上的所有点中，求它沿固定方向 $(1,0,0)$ 的方向导数的最大值及取值点。
\end{exercise}
\begin{exercise}{281}{常规·切平面截体}
求单位球面在 $P(1/\sqrt3,1/\sqrt3,1/\sqrt3)$ 处的切平面，以及该平面与三个坐标面围成四面体的体积。
\end{exercise}
\begin{exercise}{282}{基础·闭区域最值变式}
求 $f(x,y)=x+y$ 在三角形 $x\ge0,y\ge0,x+y\le1$ 上的最大值、最小值及取值点。
\end{exercise}
\begin{exercise}{283}{常规·成本恢复与约束优化}
某企业总成本满足 $C_x=20+x/2$、$C_y=6+y$、$C(0,0)=10000$，$x,y\ge0$。求 $C(x,y)$，并在 $x+y=50$ 下按实数产量求最省成本的 $x,y$。
\end{exercise}

\clearpage
\section{第10章\quad 重积分}
\subsection{二重积分的概念与性质}

\begin{exercise}{284}{基础·积分方程}
设 $D=[0,1]^2$，连续函数满足 $f(x,y)=xy+\frac12\iint_D f(u,v)\,\dd u\dd v$。令积分为 $C$，求 $C$ 及 $f(x,y)$。
\end{exercise}
\begin{exercise}{285}{基础·二重和极限}
将 $\lim_{n\to\infty}n^{-2}\sum_{i=1}^n\sum_{j=1}^n(i/n+j/n)$ 化为单位正方形上的二重积分并计算。
\end{exercise}
\begin{exercise}{286}{基础·黎曼和}
计算 $\lim_{n\to\infty}\frac{\pi}{2n^4}\sum_{i=1}^n\sum_{j=1}^n i^2\sin\frac{j\pi}{2n}$。提示：写成 $x^2\sin y$ 的二重积分。
\end{exercise}
\begin{exercise}{287}{基础·积分比较}
设 $D=\{x,y\geq0:x+y\leq1\}$。比较 $\iint_D(x+y)^2\,\dd x\dd y$ 与 $\iint_D(x+y)\,\dd x\dd y$ 的大小，说明依据。
\end{exercise}
\begin{exercise}{288}{基础·函数大小与积分}
设 $D=[0,1]^2$，$u=(x+y)/2$。按从小到大排列 $\iint_Du\,\dd\sigma,\iint_D\sqrt u\,\dd\sigma,\iint_D\sqrt[3]u\,\dd\sigma$，无需计算积分值。
\end{exercise}
\begin{exercise}{289}{基础·变量分离}
设 $D=[0,\pi]^2$，求 $\iint_D\sin^2x\sin^2y\,\dd x\dd y$ 及被积函数在 $D$ 上的平均值。
\end{exercise}
\begin{exercise}{290}{基础·积分估值}
设 $D=\{x^2+y^2\leq1\}$。利用 $1\leq1+x^2+y^2\leq2$ 给出 $\iint_D(1+x^2+y^2)\,\dd\sigma$ 的上下界。
\end{exercise}
\begin{exercise}{291}{基础·分片常值积分}
设 $D=[0,3]^2$，$f$ 在 $0\leq y\leq1$ 上为 $1$，在 $1<y\leq2$ 上为 $2$，在 $2<y\leq3$ 上为 $3$。求 $\iint_D f\,\dd\sigma$。
\end{exercise}
\begin{exercise}{292}{基础·积分平均值}
设 $D=\{0\leq y\leq x\leq1\}$。求 $xy$ 在 $D$ 上的平均值。
\end{exercise}
\begin{exercise}{293}{基础·小圆域平均值}
设 $D_r=\{x^2+y^2\leq r^2\}$。利用积分中值定理求 $\lim_{r\to0^+}\frac1{\pi r^2}\iint_{D_r}e^{x+y}\,\dd x\dd y$。
\end{exercise}
\begin{exercise}{294}{基础·连续函数积分极限}
设 $f$ 在原点附近连续且 $f(0,0)=2$。求 $\lim_{r\to0^+}r^{-2}\iint_{x^2+y^2\leq r^2}f(x,y)\,\dd x\dd y$。
\end{exercise}
\begin{exercise}{295}{基础·奇偶对称性}
计算 $\iint_{[-1,1]^2}(x^3+xy+y^2)\,\dd x\dd y$，指出哪些项由对称性消去。
\end{exercise}
\begin{exercise}{296}{常规·变量互换}
设 $D=\{x^2+y^2\leq1,x,y\geq0\}$。利用 $x,y$ 互换计算 $\iint_D\frac{e^x}{e^x+e^y}\,\dd x\dd y$。
\end{exercise}
\begin{exercise}{297}{基础·三角形积分}
设 $D=\{-1\leq x\leq y\leq1\}$。计算 $\iint_D(x+y)\,\dd x\dd y$。
\end{exercise}
\begin{exercise}{298}{基础·圆域对称性}
设 $D=\{x^2+y^2\leq1\}$。利用对称性与极坐标计算 $\iint_D(x^2+2y)\,\dd x\dd y$。
\end{exercise}
\begin{exercise}{299}{基础·非负函数积分}
设 $D=[0,1]^2$，$f$ 在 $D$ 上连续且非负。证明若 $\iint_D f\,\dd\sigma=0$，则 $f\equiv0$。提示：若某点函数值为正，利用连续性在其附近给出正的下界。
\end{exercise}
\begin{exercise}{300}{常规·积分方程小综合}
设 $D=[0,1]^2$，连续函数满足 $f(x,y)=x+y+\frac12\iint_D f(u,v)\,\dd u\dd v$。求 $f$ 并计算它在 $D$ 上的平均值。
\end{exercise}

\subsection{二重积分的计算方法}
\subsubsection{积分区域与交换积分次序}
\begin{exercise}{301}{基础·交换积分次序}
将 $\int_0^1\dd x\int_x^1 f(x,y)\,\dd y$ 交换积分次序，其中 $f$ 连续；取 $f=y$ 计算积分。
\end{exercise}
\begin{exercise}{302}{基础·半圆域换序}
将 $\int_0^1\dd y\int_{-\sqrt{1-y^2}}^{\sqrt{1-y^2}}f(x,y)\,\dd x$ 交换次序，其中 $f$ 在闭半圆上连续。
\end{exercise}
\begin{exercise}{303}{基础·分块区域合并}
设 $f$ 连续。把 $\int_0^1\dd x\int_0^x f\,\dd y+\int_1^2\dd x\int_0^{2-x}f\,\dd y$ 改写为一个先对 $x$ 积分的累次积分。
\end{exercise}
\begin{exercise}{304}{基础·三角形换序}
设 $f$ 连续。将 $\int_0^1\dd y\int_0^{2y}f\,\dd x+\int_1^2\dd y\int_0^{4-2y}f\,\dd x$ 交换次序。
\end{exercise}
\begin{exercise}{305}{基础·曲线区域面积}
设 $D=\{0\leq x\leq\pi,0\leq y\leq\sin x\}$。写出先对 $x$ 积分的面积积分并求 $D$ 的面积。
\end{exercise}
\begin{exercise}{306}{常规·抛物线域换序}
设 $f$ 连续。将 $\int_1^2\dd y\int_y^{y^2}f(x,y)\,\dd x$ 交换次序，写明需要分段的点。
\end{exercise}
\begin{exercise}{307}{基础·区域重新组合}
设 $f$ 连续。将 $\int_0^1\dd x\int_0^x f\,\dd y+\int_1^2\dd x\int_0^1 f\,\dd y$ 改为先对 $x$ 积分的累次积分。
\end{exercise}
\begin{exercise}{308}{基础·变限二重积分}
设 $f$ 连续，$F(t)=\int_0^t\dd y\int_y^t f(x)\,\dd x$，$t\geq0$。先换序，再求 $F'(t)$。
\end{exercise}
\subsubsection{直角坐标计算与分区}
\begin{exercise}{309}{基础·矩形域计算}
计算 $\iint_{[0,1]^2}e^{x-y}\,\dd x\dd y$。
\end{exercise}
\begin{exercise}{310}{基础·区域可加性}
设 $D=[0,1]^2$，$D_-=D\cap\{x+y\leq1\}$。求 $x+y$ 在 $D_-$ 及 $D$ 上的积分。
\end{exercise}
\begin{exercise}{311}{基础·指数积分}
计算 $\int_0^1\dd x\int_0^1xe^{xy}\,\dd y$。
\end{exercise}
\begin{exercise}{312}{基础·三角形域计算}
设 $D=\{1\leq y\leq x\leq2\}$。求 $\iint_Dxy\,\dd x\dd y$。
\end{exercise}
\begin{exercise}{313}{基础·曲边薄片质量}
设 $D$ 由 $y=\sqrt{x}$ 与 $y=x^2$ 围成，面密度为 $\rho=1$。求其质量并写出交点。
\end{exercise}
\begin{exercise}{314}{常规·三角形质心}
设 $D=\{x,y\geq0,x+y\leq1\}$，面密度为 $\rho=x+y$。求质量及质心坐标。
\end{exercise}
\begin{exercise}{315}{基础·三角函数积分}
设 $D=\{0\leq y\leq x\leq\pi\}$。求 $\iint_D\cos(x+y)\,\dd x\dd y$。
\end{exercise}
\begin{exercise}{316}{基础·双曲线边界}
设 $D=\{1\leq x\leq2,1/x\leq y\leq x\}$。求 $\iint_D x\,\dd x\dd y$。
\end{exercise}
\begin{exercise}{317}{基础·非均匀薄片质量}
设 $D=\{0\leq x\leq1,x\leq y\leq2x\}$，密度为 $\rho=x+y$。求质量。
\end{exercise}
\begin{exercise}{318}{基础·梯形域积分}
设 $D=\{0\leq x\leq1,0\leq y\leq1+x\}$。求 $\iint_D(1+x)\,\dd x\dd y$。
\end{exercise}
\begin{exercise}{319}{基础·多项式积分}
设 $D=\{0\leq y\leq2,y/2\leq x\leq y\}$。计算 $\iint_D(x^2+y^2)\,\dd x\dd y$。
\end{exercise}
\begin{exercise}{320}{基础·平行四边形换元}
设 $D=\{1\leq y\leq3,y-1\leq x\leq y\}$。用 $u=y-x,v=y$ 计算 $\iint_D(y-x)\,\dd x\dd y$，写出雅可比绝对值。
\end{exercise}
\begin{exercise}{321}{基础·三角形形心}
设均匀薄片的区域为顶点 $(0,0),(1,2),(2,1)$ 所成的三角形。用二重积分求形心。
\end{exercise}
\begin{exercise}{322}{基础·积分次序选择}
设 $D=\{0\leq y\leq1,0\leq x\leq y\}$。选择适当次序计算 $\iint_D\cos(y^2)\,\dd x\dd y$。
\end{exercise}
\begin{exercise}{323}{常规·可去点与积分}
令 $f(x)=\sin x/x$（$x\ne0$），$f(0)=1$。设 $D=\{0\leq x\leq1,0\leq y\leq x\}$，计算 $\iint_Df(x)\,\dd x\dd y$。
\end{exercise}
\begin{exercise}{324}{基础·换序计算}
计算 $\int_0^1\dd x\int_x^1 e^{-y^2}\,\dd y$。
\end{exercise}
\begin{exercise}{325}{基础·含高斯因子的积分}
设 $D=\{0\leq x\leq y\leq1\}$。计算 $\iint_Dx^2e^{-y^2}\,\dd x\dd y$。提示：先对 $x$ 积分，再令 $u=y^2$。
\end{exercise}
\begin{exercise}{326}{基础·区域宽度}
设 $D=\{0\leq x\leq1,0\leq y\leq x\}$。求 $\iint_De^{x^2}\,\dd x\dd y$。
\end{exercise}
\begin{exercise}{327}{基础·换序与分离}
将 $\int_0^1\dd x\int_x^1(y+xye^{y^2})\,\dd y$ 交换积分次序并计算。
\end{exercise}
\begin{exercise}{328}{基础·换序与代换}
计算 $\int_0^1\dd x\int_x^1\frac{x}{\sqrt{1+y^3}}\,\dd y$。提示：先换序。
\end{exercise}
\begin{exercise}{329}{基础·对数边界积分}
设 $D=\{1\leq x\leq e,0\leq y\leq\ln x\}$。计算 $\iint_Dx^{-1}\,\dd x\dd y$。
\end{exercise}
\begin{exercise}{330}{常规·换序与分部积分}
设 $f\in C^1[0,1]$，$f(0)=0,f(1)=1,\int_0^1 f(y)\,\dd y=1/3$。先换序，再计算 $\int_0^1\dd x\int_0^x f'(y)\,\dd y$。
\end{exercise}
\begin{exercise}{331}{常规·支集分段}
设 $f(x)=1$（$0\leq x\leq1$），其余处为 $0$。求 $H(y)=\int_{-\infty}^{+\infty}f(x)f(y-x)\,\dd x$ 的分段表达式。
\end{exercise}
\begin{exercise}{332}{基础·支集与区域}
设 $f(x,y)=xy$ 在 $[0,1]^2$ 上定义，其余处为零。求 $\iint_{x+y\leq1}f(x,y)\,\dd x\dd y$。
\end{exercise}
\begin{exercise}{333}{基础·绝对值分区}
设 $D=[0,1]^2$。沿 $y=x$ 分区计算 $\iint_D|y-x|\,\dd x\dd y$。
\end{exercise}
\begin{exercise}{334}{基础·最大值函数}
设 $D=[0,1]^2$。计算 $\iint_D\max\{x,y\}\,\dd x\dd y$，写出分区界线。
\end{exercise}
\begin{exercise}{335}{常规·对称分区积分}
设 $D=[0,1]^2$。利用 $x,y$ 互换，计算 $\iint_De^{\max\{x^2,y^2\}}\,\dd x\dd y$。
\end{exercise}
\begin{exercise}{336}{常规·根式边界换元}
设 $D=\{x,y\geq0,\sqrt x+\sqrt y\leq1\}$。用 $u=\sqrt x,v=\sqrt y$ 求 $D$ 的面积，写出换元后的积分区域。
\end{exercise}
\subsubsection{极坐标与圆形边界}
\begin{exercise}{337}{基础·扇环域}
设 $D=\{(r\cos\theta,r\sin\theta):1\leq r\leq2,\pi/6\leq\theta\leq\pi/3\}$。求 $\iint_D\arctan(y/x)\,\dd x\dd y$。
\end{exercise}
\begin{exercise}{338}{基础·圆域根式积分}
设 $D=\{x^2+y^2\leq1,x,y\geq0\}$。用极坐标计算 $\iint_D\sqrt{1-x^2-y^2}\,\dd x\dd y$。
\end{exercise}
\begin{exercise}{339}{基础·半球体积}
设 $a>0$。计算 $\iint_{x^2+y^2\leq a^2}\sqrt{a^2-x^2-y^2}\,\dd x\dd y$，说明它代表的体积。
\end{exercise}
\begin{exercise}{340}{基础·圆盘平均距离}
设 $D=\{x^2+y^2\leq1,x,y\geq0\}$。求 $\iint_D\sqrt{x^2+y^2}\,\dd x\dd y$ 及距离的平均值。
\end{exercise}
\begin{exercise}{341}{基础·平移圆域}
设 $D=\{x^2+(y-1)^2\leq1\}$。作平移换元计算 $\iint_D y\,\dd x\dd y$。
\end{exercise}
\begin{exercise}{342}{基础·环域对数积分}
设 $D=\{1\leq x^2+y^2\leq4\}$。用极坐标计算 $\iint_D\ln(x^2+y^2)\,\dd x\dd y$。
\end{exercise}
\begin{exercise}{343}{常规·对称性与极坐标}
设 $D=\{x^2+y^2\leq1,x,y\geq0\}$，原点处将被积函数定义为零。利用变量互换计算 $\iint_D\frac{x\cos(x^2+y^2)}{x+y}\,\dd x\dd y$。
\end{exercise}
\begin{exercise}{344}{基础·径向反常积分}
用圆心截断 $\varepsilon\leq\sqrt{x^2+y^2}\leq1$ 定义并计算 $\iint_{x^2+y^2\leq1}(x^2+y^2)^{-1/2}\,\dd x\dd y$。
\end{exercise}
\begin{exercise}{345}{基础·扇形径向积分}
设 $D=\{x^2+y^2\leq4,0\leq y\leq x\}$。用极坐标计算 $\iint_D\sqrt{x^2+y^2}\,\dd x\dd y$。
\end{exercise}
\begin{exercise}{346}{常规·三角形反常积分}
计算 $\int_0^1\dd x\int_x^{\sqrt3x}\frac1{\sqrt{x^2+y^2}}\,\dd y$，原点按圆心截断极限理解。提示：极坐标中 $r\cos\theta\leq1$。
\end{exercise}
\begin{exercise}{347}{基础·偏心圆域}
设 $D=\{x^2+y^2\leq2x,y\geq0\}$。写出极坐标范围并求 $D$ 的面积。
\end{exercise}
\begin{exercise}{348}{基础·半圆域对称性}
设 $D=\{x^2+y^2\leq1,y\geq0\}$。求 $\iint_D(x+y)\,\dd x\dd y$。
\end{exercise}
\begin{exercise}{349}{基础·偏心圆的矩}
设 $D=\{(x-1)^2+y^2\leq1\}$。利用平移及对称性求 $\iint_D(x+y^2)\,\dd x\dd y$。
\end{exercise}
\begin{exercise}{350}{常规·圆域积分方程}
设 $D=\{x^2+y^2\leq1\}$，连续函数满足 $f(x,y)=x^2+y^2+\frac1{2\pi}\iint_Df\,\dd\sigma$。求 $f$。
\end{exercise}
\begin{exercise}{351}{基础·径向分区}
设 $D=\{x^2+y^2\leq1\}$。计算 $\iint_D\min\{x^2+y^2,1/4\}\,\dd x\dd y$，写出分界圆。
\end{exercise}
\begin{exercise}{352}{基础·绝对值与对称性}
设 $D=\{x^2+y^2\leq1\}$。计算 $\iint_D|xy|\,\dd x\dd y$。
\end{exercise}
\subsubsection{一般换元与无界积分}
\begin{exercise}{353}{基础·线性换元}
设 $D=\{0\leq x+y\leq1,0\leq x-y\leq1\}$。令 $u=x+y,v=x-y$，计算 $\iint_D(x+y)\,\dd x\dd y$，写出雅可比绝对值。
\end{exercise}
\begin{exercise}{354}{基础·椭圆换元}
设 $D=\{x^2/4+y^2\leq1\}$。用 $x=2r\cos\theta,y=r\sin\theta$ 计算 $\iint_D(x^2/4+y^2)\,\dd x\dd y$。
\end{exercise}
\begin{exercise}{355}{常规·和与比值换元}
设 $D=\{x,y\geq0,x+y\leq1\}$，原点处赋值为 $0$。用 $u=x+y,v=x/(x+y)$ 计算 $\iint_De^{x/(x+y)}\,\dd x\dd y$；写出逆变换及雅可比绝对值。
\end{exercise}
\begin{exercise}{356}{常规·乘积换元}
设 $D=\{x,y>0,1\leq xy\leq2,1\leq y/x\leq4\}$。用 $u=xy,v=y/x$ 求 $D$ 的面积。提示：$x=\sqrt{u/v},y=\sqrt{uv}$。
\end{exercise}
\begin{exercise}{357}{基础·菱形域换元}
设 $D=\{|x|+|y|\leq1\}$。用 $u=x+y,v=x-y$ 计算 $\iint_D(x+y)^2\,\dd x\dd y$。
\end{exercise}
\begin{exercise}{358}{基础·无界区域换序}
设 $D=\{x\geq0,x^2\leq y\leq4x^2\}$。先以 $y\leq R$ 截断并交换次序，再求 $\iint_Dxe^{-y}\,\dd x\dd y$。
\end{exercise}
\begin{exercise}{359}{基础·高斯型积分}
用半径 $R$ 的圆盘截断，计算第一象限上的 $\iint e^{-(x^2+y^2)}\,\dd x\dd y$。
\end{exercise}
\begin{exercise}{360}{基础·无界径向积分}
用第一象限半径 $R$ 的圆盘递增截断，计算 $\iint_{x,y\geq0}(1+x^2+y^2)^{-2}\,\dd x\dd y$。
\end{exercise}
\begin{exercise}{361}{基础·换元区域变式}
设 $D=\{0\leq x+y\leq2,0\leq x-y\leq1\}$。写出 $u=x+y,v=x-y$ 的逆变换，并用换元计算 $D$ 的面积。
\end{exercise}
\begin{exercise}{362}{常规·反常积分小综合}
设 $D_\varepsilon=\{\varepsilon^2\leq x^2+y^2\leq1\}$。计算 $\iint_{D_\varepsilon}(x^2+y^2)^{-1}\,\dd x\dd y$，判断 $\varepsilon\to0^+$ 时积分是否有有限极限，并与幂次 $-1/2$ 的情形比较。
\end{exercise}
\subsection{三重积分}
\begin{exercise}{363}{基础·直角坐标三重积分}
设 $\Omega=\{0\leq y\leq x\leq1,0\leq z\leq1\}$。写出累次积分并计算 $\iiint_\Omega xy\,\dd x\dd y\dd z$。
\end{exercise}
\begin{exercise}{364}{基础·积分次序选择}
设 $\Omega=\{0\leq x\leq\pi/2,0\leq y\leq x,0\leq z\leq1\}$。计算 $\iiint_\Omega\sin x\,\dd V$。
\end{exercise}
\begin{exercise}{365}{基础·椭球截面}
设 $\Omega=\{x^2+y^2/4+z^2/9\leq1,z\geq0\}$。用水平截面计算 $\iiint_\Omega z\,\dd V$。
\end{exercise}
\begin{exercise}{366}{基础·抛物面围体积}
设 $\Omega=\{x^2+y^2\leq z\leq1\}$。写出柱面坐标范围并求体积。
\end{exercise}
\begin{exercise}{367}{基础·水平截层积分}
设 $\Omega=\{1\leq z\leq2,x^2+y^2\leq z\}$。求 $\iiint_\Omega(x^2+y^2)\,\dd V$。
\end{exercise}
\begin{exercise}{368}{基础·旋转曲面围体}
将 $x=0$ 平面内的抛物线 $y^2=2z$ 绕 $z$ 轴旋转，与 $z=1$ 围成 $\Omega$。写出曲面方程并求体积。
\end{exercise}
\begin{exercise}{369}{常规·球域积分}
设 $\Omega=\{x^2+y^2+z^2\leq1,x,y,z\geq0\}$。用球面坐标计算 $\iiint_\Omega xyz\,\dd V$。
\end{exercise}
\begin{exercise}{370}{基础·球锥交域}
设 $\Omega=\{x^2+y^2+z^2\leq1,z\geq\sqrt{x^2+y^2}\}$。写出球面坐标范围并求体积，极角从正 $z$ 轴起算。
\end{exercise}
\begin{exercise}{371}{基础·半球对称性}
设 $\Omega=\{x^2+y^2+z^2\leq1,z\geq0\}$。利用对称性计算 $\iiint_\Omega(x+y)\,\dd V$，再计算 $\iiint_\Omega z\,\dd V$。
\end{exercise}
\begin{exercise}{372}{基础·球域二次矩}
设 $\Omega$ 为单位球体。利用 $x,y,z$ 的对称性计算 $\iiint_\Omega(x^2+2y^2+3z^2)\,\dd V$。
\end{exercise}
\begin{exercise}{373}{基础·球体转动惯量}
半径为 $R>0$、密度为 $1$ 的实心球以原点为球心。求其关于 $z$ 轴的转动惯量。
\end{exercise}
\begin{exercise}{374}{常规·球锥域一阶矩}
设 $\Omega=\{x^2+y^2+z^2\leq1,z\geq\sqrt{x^2+y^2}\}$。计算 $\iiint_\Omega(x+z)\,\dd V$。
\end{exercise}
\begin{exercise}{375}{基础·径向函数积分}
分别计算单位球体上的 $\iiint(x^2+y^2+z^2)\,\dd V$ 与单位圆盘上的 $\iint(x^2+y^2)\,\dd\sigma$，求两者的比值。
\end{exercise}
\begin{exercise}{376}{基础·圆柱体积分}
设 $\Omega=\{x^2+y^2\leq1,0\leq z\leq2\}$。用柱面坐标计算 $\iiint_\Omega(x^2+y^2)\,\dd V$。
\end{exercise}
\begin{exercise}{377}{基础·球冠体积变式}
设 $\Omega=\{x^2+y^2+z^2\leq1,z\geq1/2\}$。用水平截面法求体积。
\end{exercise}
\begin{exercise}{378}{常规·抛物体质心小综合}
设 $\Omega=\{x^2+y^2\leq z\leq1\}$，体密度为 $\rho=1+z$。求质量及质心坐标。
\end{exercise}
\subsection{重积分的应用}
\begin{exercise}{379}{基础·平面三角形面积}
用曲面面积积分求平面 $x+y+z=1$ 在第一卦限部分的面积。
\end{exercise}
\begin{exercise}{380}{基础·圆柱侧面积}
用参数曲面的面积公式求圆柱面 $x^2+y^2=4,0\leq z\leq3$ 的侧面积。
\end{exercise}
\begin{exercise}{381}{基础·球带面积}
求单位球面上 $0\leq z\leq1/2$ 部分的面积。可用球面参数或球带公式。
\end{exercise}
\begin{exercise}{382}{常规·薄片质量与质心}
设 $D=\{0\leq x\leq1,0\leq y\leq x\}$，面密度为 $\rho=x$。求质量及质心。
\end{exercise}
\begin{exercise}{383}{基础·对称密度的质心}
设 $\Omega$ 为单位球体，体密度为 $\rho=1+x^2+y^2+z^2$。求质量，并利用对称性确定质心。
\end{exercise}
\begin{exercise}{384}{常规·抛物体形心}
设均匀物体为 $\Omega=\{x^2+y^2\leq z\leq1\}$。求其体积与形心坐标。
\end{exercise}
\begin{exercise}{385}{基础·圆锥形心}
设均匀圆锥体为 $\Omega=\{0\leq z\leq1,x^2+y^2\leq(1-z)^2\}$。利用水平截面求形心高度。
\end{exercise}
\begin{exercise}{386}{基础·椭圆薄片惯量}
设 $a,b>0$，均匀薄片 $D=\{x^2/a^2+y^2/b^2\leq1\}$ 的面密度为 $1$。求关于 $x$ 轴与 $y$ 轴的转动惯量。
\end{exercise}
\begin{exercise}{387}{基础·矩形薄片惯量}
设薄片为 $D=[-1,1]\times[0,1]$，面密度为 $1$。求关于直线 $y=1/2$ 的转动惯量。
\end{exercise}
\begin{exercise}{388}{常规·旋转体轴惯量}
设 $\Omega=\{x^2+y^2\leq z\leq1\}$，体密度为 $1$。求关于 $z$ 轴的转动惯量及惯量与质量之比。
\end{exercise}
\begin{exercise}{389}{常规·薄圆盘轴向引力}
半径为 $1$ 的均匀薄圆盘位于 $z=0$，面密度为 $1$，引力常数为 $G$。求它对 $(0,0,a)$（$a>0$）处单位质量质点的引力向量。提示：先用对称性确定方向，再用极坐标求轴向分量。
\end{exercise}
\begin{exercise}{390}{基础·上半球面积变式}
用球面参数 $x=\sin\varphi\cos\theta,y=\sin\varphi\sin\theta,z=\cos\varphi$ 求单位上半球面积，写明 $\varphi,\theta$ 的范围。
\end{exercise}
\begin{exercise}{391}{常规·圆盘非均匀密度小综合}
设 $D=\{x^2+y^2\leq1\}$，面密度为 $\rho=1+x/2$。求质量与质心坐标，并计算关于 $y$ 轴的转动惯量。
\end{exercise}

\clearpage
\section{第11章\quad 曲线积分}
\subsection{对弧长的曲线积分}
\begin{exercise}{392}{基础·常数弧长积分}
设曲线 $L$ 长度为 $a>0$。求 $\int_L3\,\dd s$，说明曲线反向是否改变结果。
\end{exercise}
\begin{exercise}{393}{基础·弧长积分非负性}
设 $f$ 在分段光滑曲线 $L$ 上连续且非负。由定义说明 $\int_Lf\,\dd s\geq0$。
\end{exercise}
\begin{exercise}{394}{基础·积分比较}
设曲线 $L$ 长度为 $2$，连续函数满足 $1\leq f\leq3$。给出 $\int_L f\,\dd s$ 的上下界。
\end{exercise}
\begin{exercise}{395}{基础·积分绝对值不等式}
设 $f$ 在分段光滑曲线 $L$ 上连续。利用 $-|f|\leq f\leq|f|$ 证明 $|\int_Lf\,\dd s|\leq\int_L|f|\,\dd s$。
\end{exercise}
\begin{exercise}{396}{基础·圆弧径向函数}
设 $L$ 为半径 $2$ 的下半圆 $x^2+y^2=4,y\leq0$。计算 $\int_L\sqrt{x^2+y^2}\,\dd s$。
\end{exercise}
\begin{exercise}{397}{基础·线段参数}
设 $L$ 连接 $(-1,0)$ 与 $(0,2)$。计算 $\int_L(2x-y+5)\,\dd s$。
\end{exercise}
\begin{exercise}{398}{基础·空间圆周}
设 $L:x^2+y^2=4,z=2$。计算 $\oint_L\arctan\frac{z^2}{x^2+y^2}\,\dd s$。
\end{exercise}
\begin{exercise}{399}{常规·大圆对称性}
设 $L:x^2+y^2+z^2=9,x+y+z=0$。利用三个坐标的对称性及球面方程求 $\oint_Lz^2\,\dd s$。
\end{exercise}
\begin{exercise}{400}{基础·椭圆约束}
设椭圆 $L:x^2/4+y^2/3=1$ 的周长为 $a$。利用对称性及曲线方程计算 $\oint_L(2xy+3x^2+4y^2)\,\dd s$，用 $a$ 表示结果。
\end{exercise}
\begin{exercise}{401}{基础·折线积分}
设 $L:y=|x|,-1\leq x\leq1$。计算 $\int_L(x+y)\,\dd s$。
\end{exercise}
\begin{exercise}{402}{基础·线段弧长平均}
设 $L$ 连接 $(1,0)$ 与 $(0,1)$。求 $\int_Lx^2\,\dd s$ 及 $x^2$ 的弧长平均值。
\end{exercise}
\begin{exercise}{403}{常规·互换对称性}
设 $L:x^2+y^2=1$。利用 $x,y$ 互换计算 $\oint_L\frac{e^x}{e^x+e^y}\,\dd s$。
\end{exercise}
\begin{exercise}{404}{基础·圆周奇函数消去}
设 $L:x^2+y^2=1,z=0$。计算 $\oint_L(xy+x^2+z)\,\dd s$。
\end{exercise}
\begin{exercise}{405}{基础·指数曲线弧长因子}
设 $L:y=e^x,0\leq x\leq1$。计算 $\int_L\sqrt{1+y^2}\,\dd s$。
\end{exercise}
\begin{exercise}{406}{基础·抛物线弧长积分}
设 $L:y^2=2x,0\leq y\leq\sqrt3$。以 $y$ 为参数计算 $\int_Ly\,\dd s$。
\end{exercise}
\begin{exercise}{407}{常规·双曲线弧积分}
设 $L$ 为 $xy=1$ 上 $1/2\leq x\leq1$ 的弧。计算 $\int_L\frac{x+y}{\sqrt{x^2+y^2}}\,\dd s$。
\end{exercise}
\begin{exercise}{408}{基础·平移圆弧参数}
设 $L:x=1+\cos t,y=\sin t,0\leq t\leq\pi/2$。求 $\int_Lxy\,\dd s$。
\end{exercise}
\begin{exercise}{409}{基础·圆心与平均值}
设 $L:(x-1)^2+(y-1)^2=1$。利用对称性求 $\oint_L(x+y)\,\dd s$。
\end{exercise}
\begin{exercise}{410}{基础·空间螺旋弧}
设 $L:x=\cos t,y=\sin t,z=t,0\leq t\leq2\pi$。计算 $\int_Lz\,\dd s$。
\end{exercise}
\begin{exercise}{411}{基础·对数曲线}
设 $L:y=\ln x,1\leq x\leq\sqrt3$。求 $\int_L\frac{e^y}{1+x^2}\,\dd s$。
\end{exercise}
\begin{exercise}{412}{常规·椭圆参数与根式}
设 $L:x^2+4y^2=4,y\geq0$。用 $x=2\cos t,y=\sin t$ 计算 $\int_L\sqrt{1+3y^2}\,\dd s$。
\end{exercise}
\begin{exercise}{413}{基础·极坐标曲线}
设 $L$ 的极坐标方程为 $r=2$，$0\leq\theta\leq\pi/2$。求 $\int_L(x-y)\,\dd s$。
\end{exercise}
\begin{exercise}{414}{基础·空间线段积分}
设 $L$ 连接 $(0,0,0)$ 与 $(1,2,2)$。计算 $\int_L(x+y+z)\,\dd s$。
\end{exercise}
\begin{exercise}{415}{基础·水平截圆}
设 $L:x^2+y^2+z^2=4,z=1$。求圆的半径并计算 $\oint_L(x^2+y^2)\,\dd s$。
\end{exercise}
\begin{exercise}{416}{常规·曲边完整边界}
设 $L$ 为 $y=x^2$ 与 $y=x$ 围成区域的完整边界。分两段计算 $\oint_Lx\,\dd s$。
\end{exercise}
\begin{exercise}{417}{基础·三角形边界}
设 $L$ 为顶点 $(0,0),(1,0),(0,1)$ 的三角形边界。分边计算 $\oint_L(x+y)\,\dd s$。
\end{exercise}
\begin{exercise}{418}{基础·圆周弧长积分变式}
设 $L:x^2+y^2=4$。计算 $\oint_L(x^2+y^2)\,\dd s$。
\end{exercise}
\begin{exercise}{419}{常规·圆形细丝质量与质心}
设 $L:x^2+y^2=1$，线密度为 $\rho=1+x/2$。求质量与质心坐标。
\end{exercise}
\subsection{对坐标的曲线积分}
\begin{exercise}{420}{基础·竖直线段}
设 $L$ 从 $(1,0)$ 到 $(1,2)$。分别计算 $\int_Ly\,\dd x$ 与 $\int_Ly\,\dd y$。
\end{exercise}
\begin{exercise}{421}{基础·常向量场}
设 $L$ 为从 $(1,2)$ 到 $(3,-1)$ 的分段光滑曲线。计算 $\int_L2\,\dd x-3\,\dd y$，并求反向积分。
\end{exercise}
\begin{exercise}{422}{基础·根式与绝对值}
设 $L:y=\sin x$，$x$ 从 $0$ 变到 $\pi$。计算 $\int_L\sqrt{1-y^2}\,\dd x$，提示根式等于 $|\cos x|$。
\end{exercise}
\begin{exercise}{423}{基础·抛物线坐标积分}
设 $L:y=x^2$，从 $(-1,1)$ 到 $(1,1)$。计算 $\int_Lx\,\dd y$。
\end{exercise}
\begin{exercise}{424}{基础·线段与方向}
设 $L$ 为从 $(3,0)$ 到 $(0,1)$ 的线段。计算 $\int_Lxy\,\dd y$。
\end{exercise}
\begin{exercise}{425}{基础·双曲线参数}
设 $L:xy=1$，从 $(1,1)$ 到 $(2,1/2)$。计算 $\int_Lx\ln y\,\dd y$。
\end{exercise}
\begin{exercise}{426}{基础·曲线约束代换}
设 $L:y=\cos x$，从 $(0,1)$ 到 $(\pi,-1)$。利用 $\dd y=-\sin x\,\dd x$ 计算 $\int_Le^y\sin x\,\dd x+(1-y)\,\dd y$。
\end{exercise}
\begin{exercise}{427}{基础·以纵坐标作参数}
设 $L:x=y^2$，从 $(1,-1)$ 到 $(1,1)$。计算 $\int_Ly\,\dd x+x\,\dd y$。
\end{exercise}
\begin{exercise}{428}{常规·摆线参数}
设 $L:x=t-\sin t,y=1-\cos t$，$t:0\to2\pi$。计算 $\int_Ly\,\dd x$。
\end{exercise}
\begin{exercise}{429}{基础·螺旋线坐标积分}
设 $L:x=2\cos t,y=2\sin t,z=t$，$t:0\to\pi$。计算 $\int_Ly\,\dd x+z\,\dd z$。
\end{exercise}
\begin{exercise}{430}{基础·全微分端点法}
设 $L$ 从 $(2,0)$ 沿椭圆 $x^2+4y^2=4$ 的第一象限弧到 $(0,1)$。利用全微分计算 $\int_Ly\,\dd x+x\,\dd y$。
\end{exercise}
\begin{exercise}{431}{基础·角形式圆周积分}
设 $L:x^2+y^2=1$ 取顺时针方向。用参数计算 $\oint_L\frac{y\,\dd x-x\,\dd y}{x^2+y^2}$。
\end{exercise}
\begin{exercise}{432}{基础·圆周约束化简}
设 $L:(x-1)^2+(y-2)^2=9$ 取逆时针方向。利用曲线方程计算 $\oint_L\frac{\dd x+\dd y}{x^2+y^2-2x-4y}$。
\end{exercise}
\begin{exercise}{433}{基础·空间全微分}
设 $L$ 从 $(1,2,3)$ 到 $(-1,4,0)$。求 $\int_Lx\,\dd x+y\,\dd y+z\,\dd z$，写出原函数。
\end{exercise}
\begin{exercise}{434}{基础·给定空间参数}
设 $L:x=\cos t,y=\sin t,z=1-\cos t$，$t:0\to2\pi$。计算 $\oint_Ly\,\dd x+z\,\dd y$。
\end{exercise}
\begin{exercise}{435}{基础·折线分段积分}
设 $L:y=1-|x|$，从 $(-1,0)$ 到 $(1,0)$。分段计算 $\int_Ly\,\dd x+x\,\dd y$。
\end{exercise}
\begin{exercise}{436}{基础·三角形边界积分}
设 $L$ 依次经过 $(0,0),(1,0),(1,1),(0,0)$。分边求 $\oint_Ly\,\dd x+x^2\,\dd y$。
\end{exercise}
\begin{exercise}{437}{基础·复合路径与原函数}
设 $L$ 从 $(0,0)$ 沿任意分段光滑路径到 $(0,2)$。找出原函数并计算 $\int_Ly\,\dd x+(x-1)\,\dd y$。
\end{exercise}
\begin{exercise}{438}{基础·曲边闭合路径}
设 $L$ 从原点沿 $y=\sqrt x$ 到 $(1,1)$，再沿 $y=x$ 回到原点。计算 $\oint_Ly\,\dd x$。
\end{exercise}
\begin{exercise}{439}{基础·空间折线积分}
设 $L$ 依次经过 $(0,0,0),(0,0,2),(1,0,2),(1,3,2)$。计算 $\int_Lz^2\,\dd x+x^2\,\dd y+y^2\,\dd z$。
\end{exercise}
\begin{exercise}{440}{常规·空间曲线平面化}
设 $L:x^2+y^2=1,z=2-y$，从正 $z$ 轴向下看为逆时针。用参数计算 $\oint_Lz\,\dd x+x\,\dd y$。
\end{exercise}
\begin{exercise}{441}{基础·乘积全微分}
设 $L$ 从 $(1,1,1)$ 到 $(2,1,3)$。利用 $\dd(xyz)$ 求 $\int_Lyz\,\dd x+zx\,\dd y+xy\,\dd z$。
\end{exercise}
\begin{exercise}{442}{基础·切向量与两类积分}
设 $L:x=\cos t,y=\sin t$，$t:0\to\pi/2$。写出同向单位切向量，将 $\int_LP\,\dd x+Q\,\dd y$ 化为对弧长的积分，其中 $P,Q$ 连续。
\end{exercise}
\begin{exercise}{443}{基础·参数反向}
设 $L:x=t,y=t^2,z=t^3$，$t:1\to0$。计算 $\int_L\dd x+2y\,\dd y+3z^2\,\dd z$。
\end{exercise}
\begin{exercise}{444}{常规·抛物线积分转换}
设 $L:y=x^2$，从 $(1,1)$ 到 $(0,0)$。计算 $\int_L\frac{\dd x+2x\,\dd y}{1+4x^2}$，并写出该积分对应的参数积分。
\end{exercise}
\begin{exercise}{445}{基础·原路返回变式}
设 $C$ 从 $(0,0)$ 沿线段到 $(1,0)$，再原路返回。分别计算 $\int_C1\,\dd s$ 与 $\int_C1\,\dd x$。
\end{exercise}
\begin{exercise}{446}{常规·两条路径小综合}
分别沿 $y=x$ 与 $y=x^2$ 从 $(0,0)$ 到 $(1,1)$ 计算 $\int_Ly\,\dd x$，据此判断该积分是否与路径无关。
\end{exercise}
\subsection{格林公式及其应用}
\subsubsection{封闭曲线、补线与奇点}
\begin{exercise}{447}{基础·闭曲线全微分}
设 $L$ 为分段光滑闭曲线。找出原函数并计算 $\oint_L(y-\cos x)\,\dd x+(x+\sin y)\,\dd y$。
\end{exercise}
\begin{exercise}{448}{基础·环流与面积}
设 $D$ 的边界分段光滑且取正向，已知 $\oint_{\partial D}(1+3y)\,\dd x+(5x-1)\,\dd y=4$。用格林公式求 $D$ 的面积。
\end{exercise}
\begin{exercise}{449}{基础·负向边界}
设 $L$ 为顶点 $(0,0),(3,0),(3,2)$ 的三角形顺时针边界。用格林公式求 $\oint_Ly\,\dd x+5x\,\dd y$。
\end{exercise}
\begin{exercise}{450}{基础·曲边区域格林公式}
设 $L$ 为 $D=\{0\leq x\leq1,x^2\leq y\leq x\}$ 的顺时针边界。计算 $\oint_Ly\,\dd x-x\,\dd y$。
\end{exercise}
\begin{exercise}{451}{基础·圆周环流}
设 $L:x^2+y^2=1$ 取逆时针方向。用格林公式求 $\oint_Lx^2y\,\dd x-xy^2\,\dd y$。
\end{exercise}
\begin{exercise}{452}{基础·扇环边界}
设 $D=\{(r\cos\theta,r\sin\theta):1\leq r\leq2,0\leq\theta\leq\pi/2\}$。计算完整正向边界上的 $\oint_{\partial D}-y\,\dd x+x\,\dd y$。
\end{exercise}
\begin{exercise}{453}{基础·抛物线弓形}
设 $L$ 为 $y=x^2$ 与 $y=1$ 围成区域的逆时针边界。用格林公式求 $\oint_L(e^x+y)\,\dd x-x\,\dd y$。
\end{exercise}
\begin{exercise}{454}{基础·菱形边界}
设 $L$ 为菱形 $|x|+|y|\leq1$ 的逆时针边界。计算 $\oint_L-y\,\dd x+2x\,\dd y$。
\end{exercise}
\begin{exercise}{455}{基础·折线补线}
设 $L$ 依次连接 $(-1,0),(0,1),(1,0)$。补上 $x$ 轴线段，用格林公式计算 $\int_Ly\,\dd x$。
\end{exercise}
\begin{exercise}{456}{常规·半圆弧补线}
设 $L:x^2+y^2=1,y\geq0$，从 $(-1,0)$ 到 $(1,0)$。用补线法计算 $\int_L(x^2-3y)\,\dd x+x\,\dd y$。
\end{exercise}
\begin{exercise}{457}{基础·曲线方程化简}
设 $L:x=y^2$，从 $(1,-1)$ 到 $(1,1)$。计算 $\int_L\frac{y\,\dd x+x\,\dd y}{y^2-x+2}$。
\end{exercise}
\begin{exercise}{458}{常规·带孔区域格林公式}
设 $D=\{1\leq x^2+y^2\leq4\}$。计算完整正向边界上的 $\oint_{\partial D}\frac{-y\,\dd x+x\,\dd y}{x^2+y^2}$，写明内外圆周的方向及格林公式的适用条件。
\end{exercise}
\begin{exercise}{459}{基础·避开奇点的圆周}
设 $L:(x-2)^2+y^2=1$ 取逆时针方向。用格林公式计算 $\oint_L\frac{-y\,\dd x+x\,\dd y}{x^2+y^2}$，说明原点是否在闭圆盘内。
\end{exercise}
\subsubsection{路径无关、原函数与全微分方程}
\begin{exercise}{460}{基础·保守场参数}
求使 $(ax+y)\,\dd x+(x+2y)\,\dd y$ 在平面内为全微分的参数 $a$，并写出一个原函数。
\end{exercise}
\begin{exercise}{461}{基础·系数匹配}
求使 $(2xy+a)\,\dd x+(bx^2+2y)\,\dd y$ 在平面内与路径无关的 $b$，其中 $a$ 为常数；求一个原函数。
\end{exercise}
\begin{exercise}{462}{基础·路径无关条件}
检验 $e^{xy}(y\,\dd x+x\,\dd y)$ 是否为全微分，并求从 $(0,0)$ 到 $(1,1)$ 的积分。
\end{exercise}
\begin{exercise}{463}{基础·分式原函数}
在半平面 $x+y>0$ 内，求 $\frac{y\,\dd x-x\,\dd y}{(x+y)^2}$ 的一个原函数。提示：考虑 $x/(x+y)$。
\end{exercise}
\begin{exercise}{464}{基础·复合函数全微分}
计算从 $(1,0)$ 到 $(2,1)$ 的 $\int_L\frac{y\,\dd x+x\,\dd y}{1+x^2y^2}$，写出原函数。
\end{exercise}
\begin{exercise}{465}{常规·反求场分量}
设 $f\in C^1(\mathbb R)$，$f(0)=1$，且 $2xy\,\dd x+f(x)\,\dd y$ 在平面内与路径无关。求 $f$ 及一个原函数。
\end{exercise}
\begin{exercise}{466}{基础·定义域条件}
在半平面 $x>0$ 内求 $\frac{-y\,\dd x+x\,\dd y}{x^2+y^2}$ 的原函数，并求从 $(1,0)$ 到 $(1,1)$ 的积分。
\end{exercise}
\begin{exercise}{467}{基础·乘积原函数}
找出 $(2xy+\cos x)\,\dd x+x^2\,\dd y$ 的原函数，求从 $(0,0)$ 到 $(1,2)$ 的积分。
\end{exercise}
\begin{exercise}{468}{基础·端点消去}
设 $L$ 从 $(0,0)$ 到 $(2,0)$。计算 $\int_L(e^x+2xy)\,\dd x+(x^2+e^{-y^2})\,\dd y$。原函数中允许保留一元定积分。
\end{exercise}
\begin{exercise}{469}{基础·多项式原函数}
计算从 $(1,2)$ 到 $(-1,1)$ 的 $\int_L(2x-y)\,\dd x+(2y-x)\,\dd y$，写出原函数。
\end{exercise}
\begin{exercise}{470}{基础·径向指数场}
设 $L$ 从 $(0,0)$ 到 $(1,1)$。利用全微分计算 $\int_Le^{x^2+y^2}(x\,\dd x+y\,\dd y)$。
\end{exercise}
\begin{exercise}{471}{基础·线性场原函数}
证明 $(x+2y)\,\dd x+(2x-y)\,\dd y$ 在平面内与路径无关，计算从 $(-1,-2)$ 到 $(3,4)$ 的积分。
\end{exercise}
\begin{exercise}{472}{基础·混合项配对}
求 $(xy^2+x)\,\dd x+(x^2y+y)\,\dd y$ 的原函数，并计算从 $(0,0)$ 到 $(1,1)$ 的积分。
\end{exercise}
\begin{exercise}{473}{基础·椭圆面积}
设 $a,b>0$，$L:x=a\cos t,y=b\sin t,t:0\to2\pi$。用 $\frac12\oint_L(x\,\dd y-y\,\dd x)$ 求椭圆面积。
\end{exercise}
\begin{exercise}{474}{常规·星形线面积}
设 $L:x=\cos^3t,y=\sin^3t,t:0\to2\pi$。用曲线积分求所围区域面积。
\end{exercise}
\begin{exercise}{475}{基础·归一化原函数}
求满足 $u(0,0)=0$ 且 $\dd u=(2x+\sin y)\,\dd x+x\cos y\,\dd y$ 的 $u(x,y)$。
\end{exercise}
\begin{exercise}{476}{基础·半平面原函数}
求 $\frac{y\,\dd x-x\,\dd y}{x^2+y^2}$ 在 $x>0$ 内满足 $u(1,0)=0$ 的原函数。
\end{exercise}
\begin{exercise}{477}{常规·参数与初值}
设 $\dd u=(x+a)e^{x+y}\,\dd x+xe^{x+y}\,\dd y$。求 $a$ 以及满足 $u(0,0)=1$ 的原函数。
\end{exercise}
\begin{exercise}{478}{基础·全微分方程}
检验 $(2x+y)\,\dd x+(x+2y)\,\dd y=0$ 为全微分方程，并求隐式通解。
\end{exercise}
\begin{exercise}{479}{基础·全微分方程初值}
求全微分方程 $(e^x+y)\,\dd x+(x+2y)\,\dd y=0$ 经过 $(0,1)$ 的积分曲线方程。
\end{exercise}
\begin{exercise}{480}{基础·原函数计算变式}
求 $2xy\,\dd x+(x^2+3y^2)\,\dd y$ 的一个原函数，并用偏导数检验。
\end{exercise}
\begin{exercise}{481}{常规·路径无关小综合}
设 $\omega=(ay+x)\,\dd x+(2x+y)\,\dd y$。求使它在整个平面内与路径无关的 $a$；在此时求从 $(0,0)$ 到 $(1,2)$ 的积分。
\end{exercise}
\subsection{曲线积分的应用}
\begin{exercise}{482}{基础·变限积分曲线弧长}
设 $L:y=\int_0^x\tan t\,\dd t,0\leq x\leq\pi/4$。求弧长。
\end{exercise}
\begin{exercise}{483}{基础·螺旋线弧长}
设 $L:x=2\cos t,y=2\sin t,z=3t,0\leq t\leq2\pi$。求弧长。
\end{exercise}
\begin{exercise}{484}{基础·抛物线细丝质量}
设 $L:y=x^2,-1\leq x\leq1$，线密度为 $\rho=|x|$。求质量及质心横坐标。
\end{exercise}
\begin{exercise}{485}{基础·空间细丝质量}
设 $L:x=t,y=t^2/2,z=t^3/3,0\leq t\leq1$，线密度为 $\rho=\sqrt{1+2y+3xz}$。求质量。
\end{exercise}
\begin{exercise}{486}{基础·细棒质心}
细棒位于 $x$ 轴的 $[0,1]$ 上，线密度为 $\rho=1+x$。求质量与质心。
\end{exercise}
\begin{exercise}{487}{基础·螺旋细丝重心}
设 $L:x=\cos t,y=\sin t,z=t,0\leq t\leq2\pi$，线密度为 $1$。求质量与重心高度。
\end{exercise}
\begin{exercise}{488}{基础·半圆细丝惯量}
设 $L:x^2+y^2=1,y\geq0$，线密度为 $1$。求关于 $x$ 轴和 $y$ 轴的转动惯量。
\end{exercise}
\begin{exercise}{489}{基础·圆形细丝轴惯量}
设 $L:x^2+y^2=4,z=0$，线密度为 $1$。求质量与关于 $z$ 轴的转动惯量。
\end{exercise}
\begin{exercise}{490}{基础·变力做功}
质点沿 $x=\cos t,y=\sin t,z=t$，$t:0\to\pi$，在力 $\bm F=(y,-x,z)$ 作用下移动。求力做的功。
\end{exercise}
\begin{exercise}{491}{基础·回复力做功}
力为 $\bm F=(-2x,-2y)$。求质点从 $(2,0)$ 沿任意分段光滑路径到 $(0,1)$ 时力做的功。
\end{exercise}
\begin{exercise}{492}{常规·径向力做功}
设 $\bm F=(x,y,z)/(x^2+y^2+z^2)$，质点沿线段从 $(1,0,0)$ 到 $(2,0,0)$。求功，写明路径上分母不为零。
\end{exercise}
\begin{exercise}{493}{基础·细丝平行轴公式}
设细丝质量为 $M>0$，质心纵坐标为 $\bar y$，线密度非负。展开 $(y-c)^2$，证明关于直线 $y=c$ 的惯量满足 $I(c)=I(\bar y)+M(c-\bar y)^2$。
\end{exercise}
\begin{exercise}{494}{常规·半圆质心与做功}
单位上半圆细丝的线密度为 $1$。求质心；再求力 $\bm F=(-y,x)$ 沿上半圆从 $(1,0)$ 逆时针到 $(-1,0)$ 所做的功。
\end{exercise}
\subsubsection{常规综合练习}
\begin{exercise}{495}{基础·边界互换对称性}
设 $L$ 为 $y=x^2$ 与 $x=y^2$ 在第一象限围成区域的完整边界。利用 $x,y$ 互换求 $\oint_L(x-y)\,\dd s$。
\end{exercise}
\begin{exercise}{496}{基础·圆周象限分解}
设 $L:x^2+y^2=1$。计算 $\oint_L(|x|+|y|)\,\dd s$。
\end{exercise}
\begin{exercise}{497}{基础·悬链线弧长因子}
设 $L:y=(e^x+e^{-x})/2,-1\leq x\leq1$。求 $\int_L|x|\,\dd s$。提示：$\sqrt{1+(y')^2}=(e^x+e^{-x})/2$。
\end{exercise}
\begin{exercise}{498}{基础·空间三角形弧长积分}
设 $L$ 为顶点 $(0,0,0),(1,0,0),(1,0,1)$ 的三角形完整边界。分边计算 $\oint_Lz\,\dd s$。
\end{exercise}
\begin{exercise}{499}{基础·扇形完整边界}
设 $L$ 为 $0\leq r\leq1,0\leq\theta\leq\pi/2$ 的扇形完整边界。分为两条半径和圆弧，计算 $\oint_L\sqrt{x^2+y^2}\,\dd s$。
\end{exercise}
\begin{exercise}{500}{基础·等值曲线积分}
设 $L$ 为 $f(x,y)=1$ 的等值线上从 $(0,1)$ 到 $(1,0)$ 的光滑弧，长度为 $\ell$，$f\in C^1$。计算 $\int_Lf\,\dd x$ 与 $\int_Lf\,\dd s$。
\end{exercise}
\begin{exercise}{501}{常规·曲边区域环流}
设 $L$ 为 $y=x^2$ 与 $y=\sqrt x$ 围成区域的顺时针边界。用格林公式计算 $\oint_L-y\,\dd x+x\,\dd y$。
\end{exercise}
\begin{exercise}{502}{基础·正弦曲线做功}
设 $L:y=\sin x$，从 $(0,0)$ 到 $(\pi,0)$。计算 $\int_Ly\,\dd x+2x\,\dd y$。
\end{exercise}
\begin{exercise}{503}{基础·空间闭合曲线}
设 $L:x=\cos t,y=\sin t,z=1$，$t:0\to2\pi$。计算 $\oint_Ly^2\,\dd x+z^2\,\dd y+x^2\,\dd z$。
\end{exercise}
\begin{exercise}{504}{基础·积分估值}
设有向曲线 $L$ 长度为 $3$，连续函数满足 $P^2+Q^2\leq4$。利用曲线积分估值不等式给出 $|\int_LP\,\dd x+Q\,\dd y|$ 的上界。
\end{exercise}
\begin{exercise}{505}{常规·线性变换与面积}
设 $L:|x|+|x+y|=1$ 取逆时针方向。用 $u=x,v=x+y$ 求包围面积，再计算 $\oint_L(x-y)\,\dd x+(x+y)\,\dd y$。
\end{exercise}
\begin{exercise}{506}{基础·含函数场的格林公式}
设 $z(x,y)=x+y$，$L$ 为逆时针单位圆周。计算 $\oint_L(z-3y)\,\dd x+(z+x^2)\,\dd y$。
\end{exercise}
\begin{exercise}{507}{基础·线性场环流}
设 $L:(x-1)^2+(y-1)^2=1$ 取逆时针方向。用格林公式求 $\oint_L-2y\,\dd x+3x\,\dd y$。
\end{exercise}
\begin{exercise}{508}{常规·余弦弧补线}
设 $L:y=\cos x$，从 $(-\pi/2,0)$ 到 $(\pi/2,0)$。用补线法计算 $\int_L(e^x-y)\,\dd x+x\,\dd y$。
\end{exercise}
\begin{exercise}{509}{基础·圆盘环流}
设 $L:x^2+y^2=4$ 取逆时针方向。用格林公式求 $\oint_L(1-x^2/4)y\,\dd x$。
\end{exercise}
\begin{exercise}{510}{基础·路径无关条件检验}
设 $f\in C^1(\mathbb R)$，$f(0)=0$，$f(x)\,\dd x+(x+y)\,\dd y$ 是否可能在平面内与路径无关？写出交叉偏导数并说明结论。
\end{exercise}
\begin{exercise}{511}{基础·同端点原函数法}
设 $L$ 为从 $(0,0)$ 到 $(1,1)$ 的任意分段光滑曲线。求 $\int_L(1-2xy)\,\dd x-x^2\,\dd y$。
\end{exercise}
\begin{exercise}{512}{常规·反求未知分量}
设 $f\in C^1(0,\infty)$，$f(1)=0$，$yf(x)\,\dd x+(x^2-2x)\,\dd y$ 在 $x>0$ 内与路径无关。求 $f$ 及一个原函数。
\end{exercise}
\begin{exercise}{513}{基础·由原函数求场分量}
已知原函数为 $u=x^2y+y^2$。求 $P,Q$ 使 $\dd u=P\,\dd x+Q\,\dd y$，并求从 $(0,0)$ 到 $(1,2)$ 的积分。
\end{exercise}
\begin{exercise}{514}{常规·偏导与端点积分}
设 $f_x=2x+y$，$f(0,y)=y^2$。求 $f$，并计算从 $(0,0)$ 到 $(1,2)$ 的 $\int_Lf_x\,\dd x+f_y\,\dd y$。
\end{exercise}
\begin{exercise}{515}{基础·三角形面积积分}
设 $L$ 依次经过 $(0,0),(2,0),(0,3),(0,0)$。计算 $\frac12\oint_L(x\,\dd y-y\,\dd x)$，说明其与三角形面积的关系。
\end{exercise}
\begin{exercise}{516}{常规·全微分方程参数}
求使 $(2xy+a)\,\dd x+(bx^2+2y)\,\dd y=0$ 为全微分方程的 $b$，其中 $a$ 为常数；求经过 $(0,0)$ 的积分曲线方程。
\end{exercise}
\begin{exercise}{517}{基础·弧长表达式}
分别写出 $y=\sin x$（$0\leq x\leq\pi$）的弧长积分，以及椭圆 $x=2\cos t,y=\sin t$（$0\leq t\leq2\pi$）的周长积分，无需求值。
\end{exercise}
\begin{exercise}{518}{基础·细棒质心变式}
设细棒位于 $x$ 轴的 $[0,2]$，线密度为 $\rho=x$。求质量与质心。
\end{exercise}
\begin{exercise}{519}{常规·旋转力场小综合}
设 $\bm F=(-y,x)$。分别求质点从 $(1,0)$ 沿单位圆上半圆逆时针到 $(-1,0)$，以及沿直径线段到同一终点时力所做的功，据此判断路径是否影响功。
\end{exercise}

\clearpage
\section{第12章\quad 曲面积分}
\subsection{第一类曲面积分}
\begin{exercise}{520}{基础·上半球对称性}
设 $S$ 为单位上半球面。利用对称性与球面参数计算 $\iint_S(x+y+z)\,\dd S$。
\end{exercise}
\begin{exercise}{521}{基础·平面面积元素}
设 $S:z=2y-1,x^2+(y-1)^2\leq1$。写出面积元素并计算 $\iint_Sz\,\dd S$。
\end{exercise}
\begin{exercise}{522}{基础·平面约束化简}
设 $S$ 为 $x/2+y/3+z/4=1$ 在第一卦限的部分。利用平面方程计算 $\iint_S(2x+4y/3+z)\,\dd S$。
\end{exercise}
\begin{exercise}{523}{基础·平面三角形积分}
设 $S=\{x+y+z=1,x,y,z\geq0\}$。计算 $\iint_Sy^2\,\dd S$。
\end{exercise}
\begin{exercise}{524}{基础·球面约束}
设 $S$ 为单位上半球面。利用球面方程计算 $\iint_Sz(x^2+y^2+z^2+1)\,\dd S$。
\end{exercise}
\begin{exercise}{525}{基础·球冠积分}
设 $S=\{x^2+y^2+z^2=1,z\geq1/2\}$。用球面参数求 $\iint_Sz\,\dd S$。
\end{exercise}
\begin{exercise}{526}{常规·分片曲面面积}
设 $S:|x|+|y|+|z|=1$。利用八个全等三角形求曲面面积及 $\iint_Sx\,\dd S$。
\end{exercise}
\begin{exercise}{527}{基础·球面二阶矩}
设 $S:x^2+y^2+z^2=R^2$，$R>0$。利用坐标互换求 $\iint_Sz^2\,\dd S$。
\end{exercise}
\begin{exercise}{528}{基础·锥面对称性}
设 $S:z=\sqrt{x^2+y^2},0\leq z\leq1$。利用对称性求 $\iint_S(xy+yz)\,\dd S$。
\end{exercise}
\begin{exercise}{529}{基础·锥面积分}
设 $S:z=\sqrt{x^2+y^2},0\leq z\leq1$。写出面积元素并计算 $\iint_Sz\,\dd S$。
\end{exercise}
\begin{exercise}{530}{基础·半球平均高度}
设 $S$ 为半径 $a>0$、球心在原点的上半球面。求 $\iint_Sz\,\dd S$ 及均匀曲面的平均高度。
\end{exercise}
\begin{exercise}{531}{常规·锥体边界分片}
设 $S$ 为 $\{\sqrt{x^2+y^2}\leq z\leq1\}$ 的完整边界。分成锥面与顶圆盘，求 $\iint_S(x^2+y^2)\,\dd S$。
\end{exercise}
\begin{exercise}{532}{基础·平面曲面多项式积分}
设 $S=\{x+y+z=1,x,y,z\geq0\}$。计算 $\iint_Sxy\,\dd S$。
\end{exercise}
\begin{exercise}{533}{常规·投影与一元换元}
设 $S=\{x+y+z=1,x,y,z\geq0\}$。计算 $\iint_S(1+x+y)^{-2}\,\dd S$。提示：投影后令 $u=x+y,v=x$。
\end{exercise}
\begin{exercise}{534}{基础·圆柱侧面积分}
设 $S:x^2+y^2=4,0\leq z\leq3$。计算 $\iint_S(x^2+y^2)^{-1}\,\dd S$。
\end{exercise}
\begin{exercise}{535}{基础·锥面参数化}
设 $S:z=\sqrt3\sqrt{x^2+y^2},0\leq z\leq\sqrt3$。写出参数和面积元素，求 $\iint_S(x^2+y^2)\,\dd S$。
\end{exercise}
\begin{exercise}{536}{基础·抛物面面积}
设 $S:z=1-x^2-y^2,z\geq0$。计算曲面面积。
\end{exercise}
\begin{exercise}{537}{常规·抛物面面积积分}
设 $S:z=x^2+y^2,x^2+y^2\leq1$。用极坐标计算 $\iint_Sz\,\dd S$。
\end{exercise}
\begin{exercise}{538}{基础·平面投影积分}
设 $S$ 为平面 $2x+2y+z=6$ 在第一卦限的部分。写出投影区域及面积元素，求 $\iint_Sx\,\dd S$。
\end{exercise}
\begin{exercise}{539}{基础·球冠对称积分}
设 $S=\{x^2+y^2+z^2=4,z\geq1\}$。计算 $\iint_S(x+y+z)\,\dd S$。
\end{exercise}
\begin{exercise}{540}{基础·球面积分变式}
设 $S$ 为单位上半球面。利用球面参数计算 $\iint_Sz^2\,\dd S$。
\end{exercise}
\begin{exercise}{541}{常规·锥面质量与重心}
设 $S:z=\sqrt{x^2+y^2},0\leq z\leq1$，面密度为 $\rho=z$。求质量与重心高度。
\end{exercise}
\subsection{第二类曲面积分}
\begin{exercise}{542}{基础·球面内侧投影}
设 $S$ 为单位球面在第一卦限的部分，取内侧。将 $\iint_Sz\,\dd x\dd y$ 化为平面二重积分并计算。
\end{exercise}
\begin{exercise}{543}{基础·平面下侧}
设 $S:z=0,x^2+y^2\leq1$，取下侧。计算 $\iint_S(x^2+y^2)\,\dd x\dd y$。
\end{exercise}
\begin{exercise}{544}{基础·下半球投影}
设 $S:x^2+y^2+z^2=1,z\leq0$，取下侧。计算 $\iint_Sx^2\,\dd x\dd y$，并说明改取上侧时的结果。
\end{exercise}
\begin{exercise}{545}{基础·上半球通量投影}
设 $S$ 为单位上半球面，取上侧。计算 $\iint_Sz^2\,\dd x\dd y$。
\end{exercise}
\begin{exercise}{546}{基础·闭柱体分面通量}
设 $S$ 为 $x^2+y^2\leq1,1\leq z\leq2$ 的完整边界，取外侧。分别求侧面及上下底面的 $\iint_Sz\,\dd x\dd y$，再求总和。
\end{exercise}
\begin{exercise}{547}{基础·四分之一柱面定向}
设 $S:x^2+y^2=1,x,y\geq0,0\leq z\leq3$，取外侧。计算 $\iint_Sx\,\dd y\dd z+y\,\dd z\dd x+z\,\dd x\dd y$。
\end{exercise}
\begin{exercise}{548}{基础·两类曲面积分转换}
设 $S=\{x+y+z=1,x,y,z\geq0\}$，取上侧。求单位法向量，将连续场 $F=(P,Q,R)$ 的通量写成第一类曲面积分。
\end{exercise}
\begin{exercise}{549}{基础·图像曲面通量}
设 $S:z=1-x^2-y^2,z\geq0$，取上侧。将 $F=(2x,2y,1)$ 的通量化为圆盘积分并计算。
\end{exercise}
\begin{exercise}{550}{基础·常法向量}
设 $S=\{x-y+z=1,x\geq0,y\leq0,z\geq0\}$，取上侧。求 $F=(1,2,1)$ 通过 $S$ 的通量。
\end{exercise}
\begin{exercise}{551}{基础·闭柱面无奇点通量}
设 $S$ 为 $x^2+y^2\leq1,-1\leq z\leq1$ 的完整边界，取外侧。分面计算 $\iint_Sx\,\dd y\dd z+z^2\,\dd x\dd y$。
\end{exercise}
\begin{exercise}{552}{基础·柱面无效分量}
设 $S:x^2+y^2=4,0\leq z\leq2$，取外侧。计算 $\iint_Sx\,\dd y\dd z+y\,\dd z\dd x+(z+1)\,\dd x\dd y$。
\end{exercise}
\begin{exercise}{553}{基础·抛物面下侧}
设 $S:z=(x^2+y^2)/2,0\leq z\leq2$，取下侧。计算 $\iint_Sz\,\dd x\dd y$。
\end{exercise}
\begin{exercise}{554}{常规·半球法向量}
设 $S$ 为单位上半球面，取外侧。利用球面单位法向量求 $F=(xz,yz,z^2)$ 的通量。
\end{exercise}
\begin{exercise}{555}{基础·柱面参数积分}
设 $S:x^2+y^2=1,0\leq z\leq2$，取外侧。用参数计算 $\iint_Sxz\,\dd y\dd z$。
\end{exercise}
\begin{exercise}{556}{基础·抛物面方向}
设 $S:z=x^2+y^2,z\leq1$，取下侧。求 $\iint_Sx^2\,\dd y\dd z+z\,\dd x\dd y$。
\end{exercise}
\begin{exercise}{557}{基础·四分之一柱面}
设 $S:x^2+y^2=1,x,y\geq0,0\leq z\leq1$，取外侧。求 $\iint_Sy\,\dd y\dd z$。
\end{exercise}
\begin{exercise}{558}{基础·斜面下侧}
设 $S=\{x+y-z=1,x,y\geq0,z\leq0\}$，取下侧。写出投影区域，计算 $\iint_S\dd y\dd z+z\,\dd x\dd y$。
\end{exercise}
\begin{exercise}{559}{基础·无散场通量}
设 $S$ 为单位圆柱体 $x^2+y^2\leq1,0\leq z\leq1$ 的边界外侧。计算场 $F=(y-z,z-x,x-y)$ 的总通量。
\end{exercise}
\begin{exercise}{560}{基础·球面径向通量}
设 $S$ 为半径 $R>0$、球心在原点的球面外侧。用单位法向量求场 $F=(x,y,z)$ 的通量。
\end{exercise}
\begin{exercise}{561}{常规·四面体分面通量}
设 $S$ 为四面体 $x,y,z\geq0,x+y+z\leq1$ 的边界外侧。求 $F=(yz,zx,xy)$ 的总通量，再由三个坐标面通量求斜面通量。
\end{exercise}
\begin{exercise}{562}{常规·球面径向奇点场}
设 $S:x^2+y^2+z^2=4$，取外侧。直接用球面法向量计算 $\iint_S\frac{x\,\dd y\dd z+y\,\dd z\dd x+z\,\dd x\dd y}{(x^2+y^2+z^2)^{3/2}}$，说明在球体内部直接用高斯公式缺少的条件。
\end{exercise}
\begin{exercise}{563}{基础·下半球奇偶性}
设 $S$ 为单位下半球面，取下侧。计算 $\iint_Sx^2\,\dd y\dd z+y^2\,\dd z\dd x+z^2\,\dd x\dd y$。
\end{exercise}
\begin{exercise}{564}{基础·补面计算柱面通量}
设 $S:x^2+y^2=1,0\leq z\leq2$，取外侧。补上两个底面，用高斯公式求 $F=(x,y,z)$ 的侧面通量。
\end{exercise}
\begin{exercise}{565}{常规·锥面通量}
设 $S:z=1-\sqrt{x^2+y^2},0\leq z\leq1$，取上侧。补上底面，求 $F=(x,y,2z)$ 的通量。
\end{exercise}
\begin{exercise}{566}{基础·立方体通量}
设 $S=\partial[0,1]^3$，取外侧。用高斯公式求 $F=(xy,y^2,z)$ 的总通量。
\end{exercise}
\begin{exercise}{567}{基础·切向场变式}
设 $S:z=\sqrt{x^2+y^2},1\leq z\leq2$，取下侧。写出一个法向量，说明 $F=(x,y,z)$ 通过 $S$ 的通量为零。
\end{exercise}
\begin{exercise}{568}{常规·抛物面通量小综合}
设 $S:z=1-x^2-y^2,z\geq0$，取上侧。求 $F=(x,y,z)$ 的通量，并用补面高斯公式核验。
\end{exercise}
\subsection{高斯公式}
\begin{exercise}{569}{基础·球体外通量}
设 $S$ 为单位球面外侧。用高斯公式求 $\iint_Sx\,\dd y\dd z+y\,\dd z\dd x+z\,\dd x\dd y$，写明使用公式的条件。
\end{exercise}
\begin{exercise}{570}{常规·锥面补面}
设 $S:z=\sqrt{x^2+y^2},0\leq z\leq1$，取下侧。补上顶圆盘，计算 $F=(x,2y,3(z-1))$ 通过 $S$ 的通量。
\end{exercise}
\begin{exercise}{571}{基础·上半球补面}
设 $S$ 为单位上半球面外侧。用补面法求 $F=(x,y,z+1)$ 的通量，明确底面方向。
\end{exercise}
\begin{exercise}{572}{基础·常散度通量}
设 $S$ 为 $0\leq z\leq1-x^2-y^2$ 的完整边界外侧。求场 $F=(y,z,x+z)$ 的总通量。
\end{exercise}
\begin{exercise}{573}{基础·半球体二次矩}
设 $S$ 为单位上半球体的完整边界外侧。用高斯公式求 $F=(x^3,y^3,z^3)$ 的总通量。
\end{exercise}
\begin{exercise}{574}{基础·全边界与侧面}
设圆柱体 $\Omega=\{x^2+y^2\leq1,0\leq z\leq3\}$。分别求场 $F=(0,0,z)$ 通过完整边界和侧面的外通量。
\end{exercise}
\begin{exercise}{575}{基础·立方体与对称性}
分别求 $F=(x^2,y^2,z^2)$ 通过 $[0,1]^3$ 与 $[-1,1]^3$ 边界外侧的通量。
\end{exercise}
\begin{exercise}{576}{基础·球面多项式通量}
设 $S$ 为半径 $2$、球心在原点的球面外侧。用高斯公式求 $F=(x^3,y^3,z^3)$ 的总通量。
\end{exercise}
\begin{exercise}{577}{基础·通量恢复体积}
设 $S$ 是围成区域 $\Omega$ 的分片光滑闭曲面，取外侧，已知 $F=(x,y,z)$ 的通量为 $12$。求 $\Omega$ 的体积。
\end{exercise}
\begin{exercise}{578}{基础·单位立方体}
设 $S=\partial[0,1]^3$，取外侧。求 $F=(4xz,-y^2,yz)$ 的通量。
\end{exercise}
\begin{exercise}{579}{常规·斜顶柱体}
设 $\Omega=\{x^2+y^2\leq1,0\leq z\leq1-x\}$。用高斯公式求场 $F=(x,y,z)$ 通过完整边界的外通量。
\end{exercise}
\begin{exercise}{580}{基础·绕轴旋转面}
设 $S:y=x^2+z^2,0\leq y\leq1$，取法向量的 $y$ 分量为负的一侧。补上 $y=1$ 的圆盘，求场 $F=(x,0,z)$ 的通量。
\end{exercise}
\begin{exercise}{581}{基础·非零补面通量}
设 $S$ 为半径 $2$、球心在原点的上半球面外侧。用补面法求场 $F=(0,yz,2)$ 的通量。
\end{exercise}
\begin{exercise}{582}{基础·方向余弦通量}
设 $S$ 为单位球面外侧，$n=(\cos\alpha,\cos\beta,\cos\gamma)$。用高斯公式求 $\iint_S(x^3\cos\alpha+y^3\cos\beta+z^3\cos\gamma)\,\dd S$。
\end{exercise}
\begin{exercise}{583}{基础·锥面方向余弦}
设 $S:z=\sqrt{x^2+y^2},0\leq z\leq1$，取下侧。写出单位法向量并计算 $F=(0,0,z^2)$ 的通量。
\end{exercise}
\begin{exercise}{584}{常规·开曲面常数项}
设 $S:z=1-x^2-y^2,z\geq0$，取上侧。用补面法求 $F=(x,y,z-1)$ 的通量。
\end{exercise}
\begin{exercise}{585}{常规·椭圆投影}
设 $S:z=1-x^2-y^2/4,z\geq0$，取上侧。求场 $F=(x,y,0)$ 的通量，明确椭圆投影换元的雅可比因子。
\end{exercise}
\begin{exercise}{586}{基础·对称项消去}
设 $S:z=x^2+y^2,z\leq1$，取上侧。用直接投影法求 $F=(x^2,y^2,z-1)$ 的通量。
\end{exercise}
\begin{exercise}{587}{基础·四面体通量}
设 $S$ 为 $x,y,z\geq0,2x+y+2z\leq2$ 的完整边界外侧。用高斯公式求 $F=(x^2+1,-2y,3z)$ 的通量。
\end{exercise}
\begin{exercise}{588}{常规·底面非零的补面法}
设 $S:z=1-x^2-y^2,z\geq0$，取上侧。用补面法求 $F=(-x,-y,x^2-z)$ 的通量。
\end{exercise}
\begin{exercise}{589}{基础·抛物体闭边界}
设 $S$ 为 $x^2+y^2\leq z\leq1$ 的完整边界外侧。用高斯公式求 $F=(x,y,z^2)$ 的通量。
\end{exercise}
\begin{exercise}{590}{基础·四面体坐标矩}
设 $S$ 为 $x,y,z\geq0,x+y+z\leq1$ 的边界外侧。求 $\iint_Sxy\,\dd y\dd z+yz\,\dd z\dd x+xz\,\dd x\dd y$。
\end{exercise}
\begin{exercise}{591}{常规·球锥区域通量}
设 $\Omega=\{x^2+y^2+z^2\leq1,z\geq\sqrt{x^2+y^2}\}$。用球面坐标求场 $F=(x,y,z)$ 通过完整边界的外通量。
\end{exercise}
\begin{exercise}{592}{基础·锥顶连续场}
设 $S:z=\sqrt{x^2+y^2},0\leq z\leq1$，取下侧。用补面法求 $F=(x,y,e^z)$ 的通量。
\end{exercise}
\begin{exercise}{593}{基础·球壳通量}
设 $S$ 为球壳 $1\leq x^2+y^2+z^2\leq4$ 的完整边界，取相对于球壳的外侧。求 $F=(x,y,z)$ 的总通量。
\end{exercise}
\begin{exercise}{594}{基础·球壳方向变式}
设 $S$ 为 $1\leq\sqrt{x^2+y^2+z^2}\leq2$ 的完整边界，取外侧。分别求常向量场 $F=(0,0,1)$ 通过内外球面的通量，并求总和。
\end{exercise}
\begin{exercise}{595}{常规·高斯公式小综合}
设 $S$ 为单位球面外侧，$F=(ax,y,z)$，$a\in\mathbb R$。求通量，并求使通量为零的 $a$；说明场的光滑性与积分面闭合性。
\end{exercise}
\subsection{斯托克斯公式}
\begin{exercise}{596}{基础·三角形旋度积分}
设 $L$ 为 $x+y+z=1$ 在第一卦限的三角形边界，与上侧法向符合右手定则。用斯托克斯公式求 $\oint_Ly\,\dd x+z\,\dd y+x\,\dd z$。
\end{exercise}
\begin{exercise}{597}{基础·竖直圆盘选面}
设 $L:y^2+z^2=1,x=0$，方向与法向量 $(1,0,0)$ 符合右手定则。计算 $\oint_L-y\,\dd z+z\,\dd y$。
\end{exercise}
\begin{exercise}{598}{基础·三角形边界方向}
设 $L$ 沿 $(1,0,0)\to(0,1,0)\to(0,0,1)\to(1,0,0)$。用斯托克斯公式计算 $\oint_Lz\,\dd x+x\,\dd y+y\,\dd z$，写出所取法向量。
\end{exercise}
\begin{exercise}{599}{基础·水平正方形环量}
设 $L$ 为 $0\leq x,y\leq1,z=1$ 的边界，从正 $z$ 轴向下看为逆时针。求 $\oint_L(y-z)\,\dd x+(z-x)\,\dd y$。
\end{exercise}
\begin{exercise}{600}{常规·球面大圆环量}
设 $L:x^2+y^2+z^2=1,x+y+z=0$，方向与法向量 $(1,1,1)$ 符合右手定则。用平面圆盘求 $\oint_Ly\,\dd x+z\,\dd y+x\,\dd z$。
\end{exercise}
\begin{exercise}{601}{基础·水平截圆}
设 $L:x^2+y^2=4,z=2$，从正 $z$ 轴向下看为逆时针。求 $\oint_L3y\,\dd x-xz\,\dd y+yz^2\,\dd z$。
\end{exercise}
\begin{exercise}{602}{基础·精确项与闭路}
设 $L:x^2+y^2=9,z=0$ 取从正 $z$ 轴向下看为逆时针的方向。求 $\oint_L2y\,\dd x+3x\,\dd y-z^2\,\dd z$。
\end{exercise}
\begin{exercise}{603}{常规·同边界换面}
设 $A=(y^2,xy,xz)$，$S$ 为单位上半球面上侧。用斯托克斯公式将旋度通量转到边界或单位圆盘，求 $\iint_S(\nabla\times A)\cdot n\,\dd S$，写明圆盘方向。
\end{exercise}
\begin{exercise}{604}{常规·五面曲面方向}
设 $A=(-y,x,0)$，$S$ 为单位立方体边界外侧去掉 $z=0$ 底面后的五个面。求 $\iint_S(\nabla\times A)\cdot n\,\dd S$，说明底部边界从正 $z$ 轴向下看的方向。
\end{exercise}
\begin{exercise}{605}{基础·常旋度三角形}
设 $L$ 沿 $(1,0,0)\to(0,1,0)\to(0,0,1)\to(1,0,0)$。求 $\oint_L(z-y)\,\dd x+(x-z)\,\dd y+(y-x)\,\dd z$。
\end{exercise}
\begin{exercise}{606}{基础·斜圆盘投影}
设 $L:x^2+y^2=1,z=1-x$，从正 $z$ 轴向下看为逆时针。用斯托克斯公式求 $\oint_L-y\,\dd x+x\,\dd y+\dd z$。
\end{exercise}
\begin{exercise}{607}{基础·空间参数环量}
设 $L:x=\cos t,y=\sin t,z=\cos t$，$t:0\to2\pi$。求 $\oint_Ly\,\dd x+z\,\dd y+x\,\dd z$。
\end{exercise}
\begin{exercise}{608}{基础·反向边界功}
设 $F=(y,z,x)$，$L$ 是 $x+y+z=1$ 第一卦限三角形的边界，方向与下侧法向符合右手定则。用斯托克斯公式求功，并给出反向后的值。
\end{exercise}
\begin{exercise}{609}{基础·同边界旋度通量}
设 $A\in C^1(\mathbb R^3)$，$S_1$ 为单位上半球面上侧，$S_2$ 为单位圆盘上侧。利用斯托克斯公式说明两曲面上的旋度通量相等。
\end{exercise}
\begin{exercise}{610}{常规·截圆环量小综合}
设 $L:x^2+y^2+z^2=4,z=1$，从正 $z$ 轴向下看为逆时针。求场 $A=(-yz,xz,0)$ 的环量，并用参数法核验符号。
\end{exercise}
\subsection{曲面积分的应用}
\begin{exercise}{611}{基础·无散场总通量}
设 $A=(yz,xz,xy)$，$S$ 为单位圆柱体 $x^2+y^2\leq1,0\leq z\leq1$ 的完整边界外侧。求散度与总通量。
\end{exercise}
\begin{exercise}{612}{基础·多项式通量}
设 $S=\partial[0,1]^3$，取外侧。求场 $A=(2x-z,x^2y,-xz^2)$ 的通量。
\end{exercise}
\begin{exercise}{613}{基础·平移球面常散度}
设 $S$ 为球心 $(3,-1,2)$、半径 $3$ 的球面外侧。求 $A=(2x+3z,-xz-y,y^2+2z)$ 的通量。
\end{exercise}
\begin{exercise}{614}{常规·半圆柱通量}
设 $S:y^2+z^2=1,z\geq0,0\leq x\leq1$，取柱体外侧。求 $A=(0,yz,z^2)$ 的通量。
\end{exercise}
\begin{exercise}{615}{基础·六面通量分配}
求 $A=(x,y,z)$ 通过单位立方体 $[0,1]^3$ 六个面的各自外通量与总通量。
\end{exercise}
\begin{exercise}{616}{基础·散度与球面对称性}
求 $A=(x^2+yz,y^2+xz,z^2+xy)$ 的散度，并利用对称性求单位球面外通量。
\end{exercise}
\begin{exercise}{617}{基础·复合函数散度}
求 $A=(e^{xy},\cos(xy),\cos(xz^2))$ 的散度及其在 $(0,1,0)$ 处的值。
\end{exercise}
\begin{exercise}{618}{基础·散度与旋度}
对 $A=(y^2,xy,xz)$ 分别求 $\nabla\cdot A$ 与 $\nabla\times A$。
\end{exercise}
\begin{exercise}{619}{基础·半圆柱体总通量}
设 $\Omega=\{0\leq x\leq1,y^2+z^2\leq1,z\geq0\}$。用高斯公式求 $A=(0,yz,z^2)$ 通过完整边界的外通量。
\end{exercise}
\begin{exercise}{620}{基础·点散度与局部通量}
设 $A=(xy^2,ye^z,x\ln(1+z^2))$，$S_r$ 为球心 $(1,1,0)$、半径 $r$ 的球面外侧。求 $\lim_{r\to0^+}\frac{\iint_{S_r}A\cdot n\,\dd S}{4\pi r^3/3}$。提示：先求散度再用积分中值定理。
\end{exercise}
\begin{exercise}{621}{基础·旋度与圆周环量}
求 $A=(x^2-y,4z,x^2)$ 的旋度，再求其沿 $z=0$ 平面内逆时针单位圆周的环量。
\end{exercise}
\begin{exercise}{622}{基础·梯度场势函数}
设 $A=(z+\sin y,x\cos y-z,x-y)$。验证其旋度为零，并求一个势函数。
\end{exercise}
\begin{exercise}{623}{基础·三角函数场旋度}
求 $A=(x^2\sin y,y^2\sin(xz),xy\sin z)$ 的旋度及其在原点的值。
\end{exercise}
\begin{exercise}{624}{基础·平面圆周环量}
设 $A=(xy,-yz,xz)$，$L$ 为 $z=0$ 平面上以 $(1,1,0)$ 为圆心、半径 $1$ 的圆周，从正 $z$ 轴向下看为逆时针。用斯托克斯公式求环量。
\end{exercise}
\begin{exercise}{625}{基础·径向函数梯度}
在 $r=\sqrt{x^2+y^2+z^2}>0$ 上令 $u=\ln r$。求 $\nabla u$ 与 $\nabla\cdot\nabla u$。
\end{exercise}
\begin{exercise}{626}{常规·径向函数拉普拉斯}
设 $u=1/\sqrt{x^2+y^2+z^2}$，定义于原点以外。直接求 $\nabla u$，并检验 $\nabla\cdot\nabla u=0$。
\end{exercise}
\begin{exercise}{627}{基础·梯度场无旋}
设 $u$ 有连续二阶偏导数。利用混合偏导数相等逐分量证明 $\nabla\times\nabla u=0$。
\end{exercise}
\begin{exercise}{628}{基础·散度旋度恒等式}
设 $A=(P,Q,R)$ 有连续二阶偏导数。逐项展开并证明 $\nabla\cdot(\nabla\times A)=0$。
\end{exercise}
\begin{exercise}{629}{常规·梯度与球面通量}
设 $u=x^2+y^2+z^2$。求 $\nabla u$，再用高斯公式求其通过半径 $2$、球心在原点的球面外侧的通量。
\end{exercise}
\subsubsection{常规综合练习}
\begin{exercise}{630}{基础·球面法向量}
设 $S:x^2+y^2+z^2=4$。求外单位法向量，并计算 $\iint_Sz^2\,\dd S$。
\end{exercise}
\begin{exercise}{631}{基础·柱面积分}
设 $S:x^2+y^2=1,0\leq z\leq1$。计算 $\iint_S(1+z^2)^{-1}\,\dd S$。
\end{exercise}
\begin{exercise}{632}{基础·球面平方函数}
设 $S$ 为单位球面。利用对称性计算 $\iint_S(x+y+1)^2\,\dd S$。
\end{exercise}
\begin{exercise}{633}{基础·平移球面积分}
设 $S:x^2+y^2+z^2=2z$。求球心、半径，再计算 $\iint_Sz\,\dd S$。
\end{exercise}
\begin{exercise}{634}{基础·球冠面积}
求单位球面位于 $z\geq\sqrt3/2$ 部分的面积。
\end{exercise}
\begin{exercise}{635}{常规·分片球面函数}
设 $S$ 为单位球面，$f=z$（$z\geq1/2$），其余处 $f=0$。计算 $\iint_Sf\,\dd S$。
\end{exercise}
\begin{exercise}{636}{基础·下半球取上侧}
设 $S$ 为单位下半球面，取上侧。计算 $\iint_S(1+z)\,\dd x\dd y$，说明该方向与球体外侧的关系。
\end{exercise}
\begin{exercise}{637}{常规·球面奇点场}
设 $S$ 为单位球面外侧，$r=\sqrt{x^2+y^2+z^2}$。直接计算 $F=(x,y,z)/r^3$ 的通量，说明原点奇点为何妨碍在整个球体上直接应用高斯公式。
\end{exercise}
\begin{exercise}{638}{基础·椭圆曲面投影}
设 $S:z=1-x^2/4-y^2,z\geq0$，取上侧。计算 $\iint_Sz\,\dd x\dd y$，写出椭圆换元。
\end{exercise}
\begin{exercise}{639}{常规·补面方向核验}
设 $S:z=x^2+y^2,z\leq1$，取上侧。用高斯公式求 $F=(x,y,z-1)$ 的通量，说明 $S$ 作为立体 $x^2+y^2\leq z\leq1$ 边界时的外侧方向。
\end{exercise}
\begin{exercise}{640}{基础·零散度参数}
设 $F=(ax,y,-3z)$。求使 $F$ 的散度恒为零的 $a$，并说明此时它通过任意分片光滑闭曲面外侧的通量。
\end{exercise}
\begin{exercise}{641}{常规·锥台侧面}
设 $S:z=\sqrt{x^2+y^2},1\leq z\leq2$，取相对于锥体的侧面外侧。补上两个圆盘，求 $F=(y,-x,z^2)$ 的通量。
\end{exercise}
\begin{exercise}{642}{基础·斜截柱面环量}
设 $L:x^2+y^2=1,x-y+z=2$，从正 $z$ 轴向下看为逆时针。用斯托克斯公式求 $\oint_L-y\,\dd x+x\,\dd y$。
\end{exercise}
\begin{exercise}{643}{基础·水平圆周环量}
设 $L:z=x^2+y^2,z=4$，从正 $z$ 轴向下看为逆时针。计算 $\oint_L2y\,\dd x+xz\,\dd y-yz^2\,\dd z$。
\end{exercise}
\begin{exercise}{644}{基础·平面约束化简}
设 $L:x^2+y^2=1,y+z=0$，从正 $z$ 轴向下看为逆时针。用 $z=-y$ 计算 $\oint_Lz\,\dd x+y\,\dd z$。
\end{exercise}
\begin{exercise}{645}{常规·斜圆盘旋度积分}
设 $L:x^2+y^2=1,z=x+y$，从正 $z$ 轴向下看为逆时针。用斯托克斯公式求 $\oint_Lxz\,\dd x+x\,\dd y+y^2\,\dd z/2$。
\end{exercise}
\begin{exercise}{646}{基础·菱形空间边界}
设 $L:x+y+z=2,|x|+|y|=1$，从正 $z$ 轴向下看为逆时针。求 $\oint_L-y\,\dd x+x\,\dd y+\dd z$。
\end{exercise}
\begin{exercise}{647}{常规·开曲线补闭合}
设 $L:x=\cos t,y=\sin t,z=0$，$t:0\to\pi$。补上直径线段，用斯托克斯公式求 $\int_L(y+z)\,\dd x+x\,\dd y+\dd z$。
\end{exercise}
\begin{exercise}{648}{基础·半球面前侧}
设 $S:x=\sqrt{1-y^2-z^2}$，取法向量 $x$ 分量非负的一侧。用投影法求 $F=(1,0,0)$ 的通量。
\end{exercise}
\begin{exercise}{649}{常规·抛物面通量应用}
设 $S:z=1+x^2+y^2,1\leq z\leq2$，取上侧。求 $F=(x,y,z)$ 的通量，写出投影区域与定向面积元素。
\end{exercise}
\begin{exercise}{650}{基础·平移球面通量}
设 $S:(x-2)^2+y^2+z^2=1$，取外侧。用高斯公式求 $F=(x,y,z)$ 的通量。
\end{exercise}
\begin{exercise}{651}{基础·曲面约束分母}
设 $S:y=x^2+z^2,0\leq y\leq1$，取法向量 $y$ 分量为负的一侧。先用曲面方程化简分母，再求 $F=(0,1,0)/(y-x^2-z^2+1)$ 的通量。
\end{exercise}
\begin{exercise}{652}{常规·旋转曲面通量}
将线段 $x=0,y=1,0\leq z\leq1$ 绕 $z$ 轴旋转得到侧面 $S$，取外侧。写出曲面方程并求 $F=(x,y,z^2)$ 的通量。
\end{exercise}
\begin{exercise}{653}{基础·锥面切向场变式}
设 $S:z=\sqrt{x^2+y^2},1\leq z\leq2$，取下侧。利用法向量求 $F=(y,-x,0)$ 的通量。
\end{exercise}
\begin{exercise}{654}{常规·斯托克斯小综合}
设 $S$ 为平面 $x+y+z=1$ 在第一卦限的部分，取上侧，其边界 $L$ 按右手定则定向。用斯托克斯公式求 $\oint_L(y-\cos z)\,\dd x+z\,\dd y+(x+x\sin z)\,\dd z$；提示：含三角函数的两项组成 $\dd(-x\cos z)$。
\end{exercise}

\clearpage
\section{第13章\quad 无穷级数}
\subsection{常数项级数的概念与性质}
\begin{exercise}{655}{基础·部分和与通项}
已知 $\sum_{k=1}^nu_k=2n/(n+2)$。求 $u_n$ 及 $\sum_{n=1}^{\infty}u_n$ 的和，注意单独核验 $n=1$。
\end{exercise}
\begin{exercise}{656}{基础·裂项求和}
求 $\sum_{n=1}^{\infty}\frac3{(3n-1)(3n+2)}$ 的前 $N$ 项和及总和。
\end{exercise}
\begin{exercise}{657}{基础·对数部分和}
求 $\sum_{n=1}^N\ln\frac{n+1}{n}$，据此判定相应无穷级数的敛散性。
\end{exercise}
\begin{exercise}{658}{基础·根式相消}
求 $\sum_{n=0}^N(\sqrt{n+1}-\sqrt n)$，据此判定无穷级数的敛散性。
\end{exercise}
\begin{exercise}{659}{基础·相邻项裂项}
求 $\sum_{n=1}^{\infty}(2^{1/(n+1)}-2^{1/n})$ 的和。
\end{exercise}
\begin{exercise}{660}{基础·三角数倒数}
利用 $1+\cdots+n=n(n+1)/2$ 求 $\sum_{n=1}^{\infty}(1+2+\cdots+n)^{-1}$ 的和。
\end{exercise}
\begin{exercise}{661}{基础·有理化裂项}
把 $\frac1{\sqrt{n+1}\sqrt n(\sqrt{n+1}+\sqrt n)}$ 化为两个相邻项之差，求从 $n=1$ 开始的级数和。
\end{exercise}
\begin{exercise}{662}{基础·等比级数参数}
求实数 $q$ 的范围，使 $\sum_{n=0}^{\infty}(2-q)^n$ 收敛，并求收敛时的和。
\end{exercise}
\begin{exercise}{663}{基础·对数公比}
求 $\sum_{n=1}^{\infty}(\ln2/2)^n$ 的和，写明公比绝对值小于 $1$。
\end{exercise}
\begin{exercise}{664}{基础·三角函数公比}
分别判定 $\sum_{n=0}^{\infty}\tan^n(\pi/6)$ 与 $\sum_{n=0}^{\infty}\tan^n(\pi/4)$ 的敛散性，收敛者求和。
\end{exercise}
\begin{exercise}{665}{基础·负公比}
求 $\sum_{n=1}^{\infty}\cos(n\pi)/3^n$ 的和，判断是否绝对收敛。
\end{exercise}
\begin{exercise}{666}{基础·正项级数幂次判别}
求使 $\sum_{n=1}^{\infty}n^{-2p}$ 与 $\sum_{n=1}^{\infty}n^{-(3-p)}$ 同时收敛的 $p$。
\end{exercise}
\begin{exercise}{667}{基础·收敛必要条件}
证明级数 $\sum_{n=1}^{\infty}u_n$ 收敛时必有 $u_n\to0$。提示：$u_n=S_n-S_{n-1}$。
\end{exercise}
\begin{exercise}{668}{基础·有限项与常数倍}
已知 $\sum_{n=1}^{\infty}u_n=S$。求 $\sum_{n=1}^{\infty}3u_n$ 与 $\sum_{n=2}^{\infty}u_n$ 的和。
\end{exercise}
\begin{exercise}{669}{基础·通项必要条件}
分别判定 $\sum_{n=1}^{\infty}\cos(1/n^2)$ 与 $\sum_{n=1}^{\infty}(-1)^ne^{-n}$ 的敛散性，给出依据。
\end{exercise}
\begin{exercise}{670}{基础·混合裂项求和}
求 $\sum_{n=1}^{\infty}(3^{-n}+1/(n(n+1)))$ 的和。
\end{exercise}
\begin{exercise}{671}{基础·通项条件确定参数}
求使 $\sum_{n=1}^{\infty}(an/(6n+1)+(-1)^n/6^n)$ 收敛的 $a$，并在收敛时求和。
\end{exercise}
\begin{exercise}{672}{基础·级数线性运算}
已知 $\sum_{n=1}^{\infty}u_n=2$。求 $\sum_{n=1}^{\infty}(u_n+3/2^n)$ 的和。
\end{exercise}
\begin{exercise}{673}{基础·正项幂次判定}
求使 $\sum_{n=1}^{\infty}(n^{1-p}+3\cdot2^{-n})$ 收敛的实数 $p$，说明临界值的情况。
\end{exercise}
\begin{exercise}{674}{基础·部分和线性组合}
已知 $\sum_{n=1}^{\infty}u_n=2$，$\sum_{k=1}^nv_k=2n/(3n+1)$。求 $\sum_{n=1}^{\infty}(2u_n-3v_n)$ 的和。
\end{exercise}
\begin{exercise}{675}{常规·奇偶项分离}
已知 $\sum_{n=1}^{\infty}n^{-2}=\pi^2/6$。求 $\sum_{n=1}^{\infty}(2n-1)^{-2}$ 与 $\sum_{n=1}^{\infty}(-1)^{n-1}n^{-2}$ 的和。
\end{exercise}
\begin{exercise}{676}{基础·收敛级数之和}
设 $\sum u_n$ 与 $\sum v_n$ 收敛。利用部分和证明 $\sum(u_n+v_n)$ 收敛且其和等于两级数和之和。
\end{exercise}
\begin{exercise}{677}{基础·抵消后的级数}
设 $u_n=1/n$，$v_n=-1/n+1/n^2$。判定 $\sum_{n=1}^{\infty}(u_n+v_n)$ 的敛散性。
\end{exercise}
\begin{exercise}{678}{基础·裂项求和变式}
求 $\sum_{n=1}^{\infty}\frac1{(n+1)(n+2)}$ 的和。
\end{exercise}
\begin{exercise}{679}{常规·裂项与余项小综合}
求 $\sum_{n=1}^{\infty}\frac1{n(n+1)}$ 的和及前 $N$ 项余项，求使余项小于 $10^{-2}$ 的最小正整数 $N$。
\end{exercise}
\subsection{正项级数}
\begin{exercise}{680}{基础·几何比较}
利用 $0<3+(-1)^n\leq4$ 判定 $\sum_{n=1}^{\infty}(3+(-1)^n)/5^n$ 的敛散性。
\end{exercise}
\begin{exercise}{681}{基础·比值判别}
用比值判别法判定 $\sum_{n=1}^{\infty}3^n/n!$ 的敛散性。
\end{exercise}
\begin{exercise}{682}{基础·极限比较}
分别判定 $\sum_{n=1}^{\infty}1/(3n-1)$ 与 $\sum_{n=1}^{\infty}(n+3)/(n(n+1)(n+2))$ 的敛散性，写出比较对象。
\end{exercise}
\begin{exercise}{683}{基础·有理式比较}
用极限比较法判定 $\sum_{n=1}^{\infty}(3n^2+4)/(n^4+1)$ 的敛散性。
\end{exercise}
\begin{exercise}{684}{基础·三角函数通项}
比较 $\sum_{n=1}^{\infty}\sin(\pi/(6n))$ 与 $\sum_{n=1}^{\infty}(\sin(\pi/6))^n$ 的敛散性。
\end{exercise}
\begin{exercise}{685}{基础·对数等价}
判定 $\sum_{n=1}^{\infty}\ln(1+n^{-3})$ 与 $\sum_{n=1}^{\infty}\ln(1+n^{-1/2})$ 的敛散性。
\end{exercise}
\begin{exercise}{686}{基础·三角函数比较}
用极限比较法判定 $\sum_{n=1}^{\infty}n^{-1}\tan(1/n)$ 的敛散性。
\end{exercise}
\begin{exercise}{687}{基础·已知上界}
设 $0\leq u_n\leq n^{-3}$。用比较判别法证明 $\sum_{n=1}^{\infty}(u_n+3^{-n})$ 收敛。
\end{exercise}
\begin{exercise}{688}{基础·渐近等价}
设 $u_n>0$ 且 $n^2u_n\to3$。分别判定 $\sum_{n=1}^{\infty}u_n$ 与 $\sum_{n=1}^{\infty}nu_n$ 的敛散性。
\end{exercise}
\begin{exercise}{689}{常规·平方改变敛散性}
设 $u_n=1-\cos(1/\sqrt n)$。用等价无穷小判定 $\sum_{n=1}^{\infty}u_n$ 与 $\sum_{n=1}^{\infty}u_n^2$ 的敛散性。
\end{exercise}
\begin{exercise}{690}{基础·阶乘比值}
用比值法分别判定 $\sum_{n=1}^{\infty}1/(2n+1)!$ 与 $\sum_{n=1}^{\infty}n!/2^n$ 的敛散性。
\end{exercise}
\begin{exercise}{691}{基础·指数与阶乘}
用比值法判定 $\sum_{n=1}^{\infty}2^n/n^2$ 与 $\sum_{n=1}^{\infty}2^nn!/n^n$ 的敛散性。
\end{exercise}
\begin{exercise}{692}{基础·奇数连乘}
用比值法判定 $\sum_{n=1}^{\infty}\frac{1\cdot3\cdots(2n-1)}{5^nn!}$ 的敛散性。
\end{exercise}
\begin{exercise}{693}{基础·指数衰减比较}
利用几何级数判定 $\sum_{n=1}^{\infty}1/((n+1)e^n)$ 的敛散性。
\end{exercise}
\begin{exercise}{694}{基础·指数主导项}
用比值法或极限比较法判定 $\sum_{n=1}^{\infty}5^n/(7^n-2^n)$ 的敛散性。
\end{exercise}
\begin{exercise}{695}{基础·多项式乘几何项}
用比值法判定 $\sum_{n=1}^{\infty}(2n+1)/2^n$ 的敛散性。
\end{exercise}
\begin{exercise}{696}{基础·有界振荡项}
利用 $0\leq\sin^2(n\pi/3)\leq1$ 判定 $\sum_{n=1}^{\infty}(3n+\sin^2(n\pi/3))/3^n$ 的敛散性，无需求和。
\end{exercise}
\begin{exercise}{697}{基础·递推项比值}
设 $a_1=1,a_{n+1}=\frac13(1+1/n)^n a_n$。用比值法判定 $\sum_{n=1}^{\infty}a_n$ 的敛散性。
\end{exercise}
\begin{exercise}{698}{基础·混合正项级数}
分别判定两部分，进而判定 $\sum_{n=1}^{\infty}(n^{-4}+4^n/(n+1)!)$ 的敛散性。
\end{exercise}
\begin{exercise}{699}{常规·阶乘商}
用比值法判定 $\sum_{n=1}^{\infty}(n!)^2/(2n+1)!$ 的敛散性，并由此求通项极限。
\end{exercise}
\begin{exercise}{700}{基础·根值判别}
用根值法判定 $\sum_{n=1}^{\infty}(2n/(4n+1))^n$ 的敛散性。
\end{exercise}
\begin{exercise}{701}{基础·指数项通项条件}
判定 $\sum_{n=1}^{\infty}3^n/(2^n(3n+1))$ 的敛散性，说明通项是否趋于零。
\end{exercise}
\begin{exercise}{702}{基础·根值法两例}
用根值法判定 $\sum_{n=1}^{\infty}((4n-3)/(3n+2))^{2n}$ 与 $\sum_{n=2}^{\infty}(1-1/n)^{n^2}$ 的敛散性。
\end{exercise}
\begin{exercise}{703}{基础·有界底数幂}
用根值法判定 $\sum_{n=1}^{\infty}(\arctan n)^n/2^n$ 的敛散性。
\end{exercise}
\begin{exercise}{704}{常规·根值法临界值}
判定 $\sum_{n=1}^{\infty}(n/(n+1))^n$ 的敛散性；若根值极限为 $1$，继续检查通项极限。
\end{exercise}
\begin{exercise}{705}{基础·比较法变式}
用极限比较法判定 $\sum_{n=1}^{\infty}1/(n^2+3n+1)$ 的敛散性。
\end{exercise}
\begin{exercise}{706}{常规·相消后比较小综合}
用 $\ln(1+t)=t-t^2/2+o(t^2)$ 判定 $\sum_{n=1}^{\infty}(1/n-\ln(1+1/n))$ 的敛散性，写出比较极限。
\end{exercise}
\subsection{任意项级数}
\begin{exercise}{707}{基础·绝对收敛}
判定 $\sum_{n=1}^{\infty}(-1)^n/(2n)!$ 的收敛类型。
\end{exercise}
\begin{exercise}{708}{基础·交错与通项条件}
分别判定 $\sum_{n=1}^{\infty}(-1)^n\sin(1/\sqrt n)$ 与 $\sum_{n=1}^{\infty}(-1)^n\cos(1/\sqrt n)$ 的收敛类型或发散，写出依据。
\end{exercise}
\begin{exercise}{709}{基础·对数交错级数}
判定 $\sum_{n=1}^{\infty}(-1)^n\ln(1+1/n)$ 的收敛类型，验证绝对值通项单调趋零。
\end{exercise}
\begin{exercise}{710}{基础·绝对收敛与拆分}
判定 $\sum_{n=1}^{\infty}(-1)^n(2+(-1)^n)/4^n$ 是否绝对收敛，并拆成几何级数求和。
\end{exercise}
\begin{exercise}{711}{常规·成对分组}
按所给顺序判定 $1-1/2+1/2-1/3+1/3-1/4+\cdots$ 的敛散性并求和。提示：先求前 $2N$ 项和，再检查奇数项部分和。
\end{exercise}
\begin{exercise}{712}{基础·交错幂次}
分别判定 $\sum_{n=1}^{\infty}(-1)^n/\sqrt n$ 与 $\sum_{n=1}^{\infty}(-1)^n/n^{3/2}$ 的收敛类型。
\end{exercise}
\begin{exercise}{713}{常规·周期符号}
把 $\sum_{n=1}^{\infty}\sin(n\pi/2)/n$ 去掉零项后写成通常的交错级数，判定其收敛类型。
\end{exercise}
\begin{exercise}{714}{常规·最终单调性}
判定 $\sum_{n=1}^{\infty}(-1)^n/(n-\ln n)$ 的收敛类型。提示：考察 $x-\ln x$ 在 $x>1$ 上的单调性。
\end{exercise}
\begin{exercise}{715}{基础·绝对值比较}
设 $|u_n|\leq1/n$。用比较法证明 $\sum_{n=1}^{\infty}u_n/n$ 绝对收敛。
\end{exercise}
\begin{exercise}{716}{基础·交错余项估计}
对 $\sum_{n=1}^{\infty}(-1)^{n-1}/(n+1)$ 判定收敛类型，并给出取前 $N$ 项后的莱布尼茨余项界。
\end{exercise}
\begin{exercise}{717}{基础·有界振荡系数}
用绝对值比较判定 $\sum_{n=1}^{\infty}(-1)^n\cos n/\pi^{n+1}$ 的收敛类型。
\end{exercise}
\begin{exercise}{718}{常规·谐波项分离}
求使 $\sum_{n=1}^{\infty}(((-1)^n+kn)/n^2)$ 收敛的 $k$，判断收敛时是否绝对收敛。
\end{exercise}
\begin{exercise}{719}{基础·交错与几何型}
判定 $\sum_{n=1}^{\infty}(-1)^{n-1}n^{-2/3}$ 与 $\sum_{n=1}^{\infty}(-1)^{n-1}(2/3)^n$ 的收敛类型。
\end{exercise}
\begin{exercise}{720}{基础·有界系数比较}
用绝对值比较判定 $\sum_{n=1}^{\infty}\frac{(-1)^n(3+(-1)^n)}{(3n+1)\sqrt{n+3}}$ 的收敛类型。
\end{exercise}
\begin{exercise}{721}{基础·交错符号与必要条件}
分别判定 $\sum_{n=1}^{\infty}(-1)^n(n/(n+1))^n$ 与 $\sum_{n=1}^{\infty}(-1)^n3^n\sin(\pi/5^n)$ 的敛散性，收敛者判断是否绝对收敛。
\end{exercise}
\begin{exercise}{722}{常规·绝对值与比值法}
用比较法及比值法判定 $\sum_{n=1}^{\infty}n^n\cos n/(3^nn!)$ 的收敛类型。
\end{exercise}
\begin{exercise}{723}{常规·平方可和}
设 $\sum_{n=1}^{\infty}u_n^2$ 收敛。利用 $2|u_n|/n\leq u_n^2+1/n^2$ 证明 $\sum_{n=1}^{\infty}u_n/n$ 绝对收敛。
\end{exercise}
\begin{exercise}{724}{常规·交错级数与平方级数}
设 $u_n=(-1)^n\tan(1/\sqrt n)$。判定 $\sum_{n=1}^{\infty}u_n$ 的收敛类型及 $\sum_{n=1}^{\infty}u_n^2$ 的敛散性。
\end{exercise}
\begin{exercise}{725}{基础·条件收敛拆分}
判定 $\sum_{n=1}^{\infty}(n+1)(-1)^n/n^2$ 的收敛类型；可用已知的交错调和级数和交错平方倒数级数作比较。
\end{exercise}
\begin{exercise}{726}{基础·绝对收敛比较}
判定 $\sum_{n=1}^{\infty}\sin n/n^2$ 的收敛类型，写出绝对值上界。
\end{exercise}
\begin{exercise}{727}{基础·莱布尼茨判别变式}
判定 $\sum_{n=2}^{\infty}(-1)^n/\ln n$ 的收敛类型，验证判别法条件。
\end{exercise}
\begin{exercise}{728}{常规·交错加正项小综合}
判定 $\sum_{n=1}^{\infty}((-1)^{n-1}/\sqrt n+1/n^2)$ 的收敛类型。提示：利用通项绝对值不小于 $1/\sqrt n-1/n^2$ 检查绝对收敛性。
\end{exercise}
\subsection{幂级数}
\begin{exercise}{729}{基础·半径与端点}
求 $\sum_{n=0}^{\infty}(3n+1)x^n/3^n$ 的收敛半径及收敛域，分别检查两个端点。
\end{exercise}
\begin{exercise}{730}{常规·两种指数系数}
求 $\sum_{n=1}^{\infty}(3^n+2^n)x^n/n^2$ 的收敛半径及收敛域，判断两个端点是否绝对收敛。
\end{exercise}
\begin{exercise}{731}{基础·稀疏幂次}
求 $\sum_{n=1}^{\infty}4^nx^{2n+1}$ 的收敛半径及收敛域。
\end{exercise}
\begin{exercise}{732}{基础·平移交错幂级数}
求 $\sum_{n=1}^{\infty}(-1)^{n-1}(x+1)^n/n$ 的收敛半径及收敛域，检查两个端点。
\end{exercise}
\begin{exercise}{733}{基础·阶乘系数}
求 $\sum_{n=0}^{\infty}x^n/(2n)!$ 的收敛半径及收敛域，无需求和。
\end{exercise}
\begin{exercise}{734}{基础·阶乘与指数系数}
求 $\sum_{n=0}^{\infty}n!x^n/3^n$ 的收敛半径及收敛域。
\end{exercise}
\begin{exercise}{735}{基础·端点条件收敛}
求 $\sum_{n=1}^{\infty}(-1)^nx^n/(2^n\sqrt n)$ 的收敛域，明确两个端点的情况。
\end{exercise}
\begin{exercise}{736}{基础·平移中心}
求 $\sum_{n=1}^{\infty}(x-2)^n/(n^2+1)$ 的收敛半径及收敛域，检查两个端点。
\end{exercise}
\begin{exercise}{737}{基础·内点绝对收敛}
已知 $\sum_{n=0}^{\infty}a_nx^n$ 在 $x=-3$ 收敛。说明它在 $x=2$ 是否绝对收敛，写出所用幂级数性质。
\end{exercise}
\begin{exercise}{738}{基础·中心与距离}
已知 $\sum_{n=0}^{\infty}a_n(x-1)^n$ 的收敛半径为 $1$。判断它在 $x=1/2$ 与 $x=3$ 处的敛散性。
\end{exercise}
\begin{exercise}{739}{基础·端点直接检验}
求 $\sum_{n=1}^{\infty}(-1)^{n-1}(x+1/2)^n/n$ 的收敛域，分别检查 $x=-3/2,1/2$。
\end{exercise}
\begin{exercise}{740}{基础·已知半径判断点}
已知 $\sum_{n=0}^{\infty}a_n(x-2)^n$ 的收敛半径为 $3$。判断它在 $x=-1/2$ 与 $x=6$ 处的敛散性。
\end{exercise}
\begin{exercise}{741}{基础·导数级数半径}
已知 $\sum_{n=0}^{\infty}a_nx^n$ 的收敛半径为 $3$。求其逐项求导后的收敛半径，并写出保证逐项求导成立的开区间。
\end{exercise}
\begin{exercise}{742}{基础·几何幂级数求和}
求 $\sum_{n=1}^{\infty}x^{n+1}/3^n$ 的收敛域及和函数，检查两个端点。
\end{exercise}
\begin{exercise}{743}{常规·望远镜幂级数}
求 $\sum_{n=1}^N(x^n/n-x^{n+1}/(n+1))$，据此求相应无穷级数在 $[-1,1]$ 上的和。
\end{exercise}
\begin{exercise}{744}{基础·对数和函数}
求 $\sum_{n=1}^{\infty}5^nx^n/n$ 的收敛域及和函数，检查两个端点。
\end{exercise}
\begin{exercise}{745}{基础·幂次平移求和}
求 $\sum_{n=1}^{\infty}(-1)^{n-1}x^{n+1}/n$ 的收敛域及和函数，检查两个端点。
\end{exercise}
\begin{exercise}{746}{常规·反正切和函数}
求 $S(x)=\sum_{n=0}^{\infty}(-1)^nx^{2n}/(2n+1)$ 的收敛域及和函数，单独写明 $S(0)$，检查两个端点。
\end{exercise}
\begin{exercise}{747}{基础·几何级数求导}
从几何级数出发求 $\sum_{n=1}^{\infty}nx^{2n-1}$ 的收敛域及和函数，检查两个端点。
\end{exercise}
\begin{exercise}{748}{基础·平移级数求和}
求 $\sum_{n=1}^{\infty}n(x+2)^n$ 的收敛域及和函数，检查两个端点。
\end{exercise}
\begin{exercise}{749}{基础·平方倒数系数变式}
求 $\sum_{n=1}^{\infty}x^n/n^2$ 的收敛半径及收敛域，检查两个端点。
\end{exercise}
\begin{exercise}{750}{常规·求导与端点小综合}
对 $F(x)=\sum_{n=1}^{\infty}(x-1)^n/n^2$，分别求原级数与逐项求导后的级数的收敛域，检查各自两个端点。
\end{exercise}
\subsection{函数展开成幂级数}
\begin{exercise}{751}{基础·复合指数展开}
将 $e^{-x^2}$ 展开为 $x$ 的幂级数，写出前四个非零项及收敛域。
\end{exercise}
\begin{exercise}{752}{基础·正弦展开}
写出 $\sin2x$ 的麦克劳林级数及收敛域，求 $x^5$ 的系数。
\end{exercise}
\begin{exercise}{753}{基础·几何展开}
将 $1/(3-x)$ 展开成 $x$ 的幂级数，写出通项及收敛域，检查两个端点。
\end{exercise}
\begin{exercise}{754}{常规·倍角公式展开}
利用 $\sin^2x=(1-\cos2x)/2$ 写出其麦克劳林级数的前四个非零项及通项，注明收敛域。
\end{exercise}
\begin{exercise}{755}{基础·识别正弦级数}
求 $\sum_{n=0}^{\infty}(-1)^nx^{2n+1}/(2n+1)!$ 的和函数及收敛域。
\end{exercise}
\begin{exercise}{756}{基础·对数展开}
由几何级数逐项积分推导 $\ln(1+x)$ 的麦克劳林展开，写出收敛域并检查两个端点。
\end{exercise}
\begin{exercise}{757}{基础·反正切展开}
由 $1/(1+x^2)$ 的几何级数推导 $\arctan x$ 的麦克劳林展开，写出收敛域并检查两个端点。
\end{exercise}
\begin{exercise}{758}{基础·有理函数展开}
将 $1/(3-x)^2$ 展开成 $x$ 的幂级数，写出通项及收敛域，检查两个端点。
\end{exercise}
\begin{exercise}{759}{常规·部分分式展开}
将 $1/((1-x)(2-x))$ 展开成 $x$ 的幂级数，写出通项及收敛域，检查两个端点。
\end{exercise}
\begin{exercise}{760}{基础·平移指数展开}
将 $e^{3x}$ 展开成 $(x-1)$ 的幂级数，写出前三项及收敛域。
\end{exercise}
\begin{exercise}{761}{基础·平移几何展开}
将 $1/(3-x)$ 展开成 $(x+1)$ 的幂级数，求收敛域并检查两个端点。
\end{exercise}
\begin{exercise}{762}{基础·对数平移展开}
将 $\ln x$ 展开成 $(x-2)$ 的幂级数，求收敛域，分别检查 $x=0,4$。
\end{exercise}
\begin{exercise}{763}{基础·负中心展开}
将 $x^{-2}$ 展开成 $(x+1)$ 的幂级数，写出通项及收敛域，检查两个端点。
\end{exercise}
\begin{exercise}{764}{常规·部分分式中心平移}
将 $1/(x(x+1))$ 展开成 $(x-1)$ 的幂级数，写出通项及收敛域，检查两个端点。
\end{exercise}
\begin{exercise}{765}{基础·指数展开变式}
写出 $e^{2x}$ 的麦克劳林级数及收敛域，并求 $x^4$ 的系数。
\end{exercise}
\begin{exercise}{766}{常规·有理函数展开小综合}
将 $x/(1-x)^2$ 展开成 $x$ 的幂级数，写出收敛域并检查端点；再求其在原点的四阶导数。
\end{exercise}
\subsection{幂级数的应用}
\begin{exercise}{767}{基础·系数求导数}
设 $f(x)=2x^2/(1+x)$。利用原点附近的幂级数求 $f^{(5)}(0)$。
\end{exercise}
\begin{exercise}{768}{常规·乘积有限阶导数}
设 $f(x)=(1+x)\arctan x$。展开到 $x^5$，求 $f^{(4)}(0)$ 与 $f^{(5)}(0)$。
\end{exercise}
\begin{exercise}{769}{基础·指数乘积导数}
设 $f(x)=xe^{-x}$。用麦克劳林系数求 $f^{(4)}(0)$。
\end{exercise}
\begin{exercise}{770}{基础·平移对数导数}
设 $f(x)=\ln(x^2+4x+5)$。利用 $x^2+4x+5=1+(x+2)^2$ 的展开求 $f^{(4)}(-2)$。
\end{exercise}
\begin{exercise}{771}{基础·平移分式导数}
将 $f(x)=(x-1)/(2x+1)$ 在 $x=1$ 附近展开到 $(x-1)^3$，求 $f'''(1)$。
\end{exercise}
\begin{exercise}{772}{常规·逐项积分求和}
由 $-\ln(1-t)=\sum_{n=1}^{\infty}t^n/n$（$|t|<1$）求 $\sum_{n=1}^{\infty}1/(2^n(n+1))$ 的和。
\end{exercise}
\begin{exercise}{773}{基础·线性加权几何和}
由几何级数逐项求导，求 $\sum_{n=0}^{\infty}(n+1)/3^n$ 的和。
\end{exercise}
\begin{exercise}{774}{常规·逐项积分和函数}
在 $|x|<1$ 内求 $F(x)=\sum_{n=0}^{\infty}(-1)^nx^{2n+1}/(2n+1)$ 的和函数，再计算 $F(1/2)$。
\end{exercise}
\begin{exercise}{775}{常规·三角级数加权}
利用 $\cos x$ 的展开及其导数求 $\sum_{n=0}^{\infty}(-1)^n(2n+1)x^{2n}/(2n)!$ 的和函数。
\end{exercise}
\begin{exercise}{776}{常规·二次加权几何和}
由几何级数求导，求 $\sum_{n=1}^{\infty}n(n+1)/2^n$ 的和。
\end{exercise}
\begin{exercise}{777}{基础·对数数项和变式}
由 $\ln(1+x)$ 的幂级数求 $\sum_{n=1}^{\infty}(-1)^{n-1}/(n2^n)$ 的和。
\end{exercise}
\begin{exercise}{778}{常规·系数与求和小综合}
写出 $\ln(1+x^2)$ 在原点的幂级数，求其四阶原点导数，并利用该展开求 $\sum_{n=1}^{\infty}(-1)^{n-1}/(n4^n)$ 的和。
\end{exercise}
\subsection{傅里叶级数}
\begin{exercise}{779}{基础·傅里叶常数项}
设 $f$ 以 $2\pi$ 为周期，在 $0<x\leq2\pi$ 上为 $f(x)=x$。按 $f\sim a_0/2+\sum_{n\geq1}(a_n\cos nx+b_n\sin nx)$ 的约定，求 $a_0$ 及级数在 $x=0$ 处的和。
\end{exercise}
\begin{exercise}{780}{基础·周期归约与跳跃平均}
将 $f(x)=x+1$（$-\pi\leq x<\pi$）作 $2\pi$ 周期延拓。求其傅里叶级数在 $x=7\pi/2,\pi,0$ 处的和，无需计算系数。
\end{exercise}
\begin{exercise}{781}{基础·分段函数的傅里叶和}
设 $2\pi$ 周期函数在 $-\pi\leq x<0$ 时为 $f(x)=x$，在 $0\leq x<\pi$ 时为零。求傅里叶级数在 $x=-\pi/2,0,\pi$ 处的和。
\end{exercise}
\begin{exercise}{782}{基础·正弦级数奇延拓}
设 $f(x)=1$（$0<x<\pi$），$b_n=\frac2\pi\int_0^\pi f(x)\sin(nx)\,\dd x$。求 $S(x)=\sum_{n=1}^{\infty}b_n\sin(nx)$ 在 $x=0,\pi/2,\pi$ 处的和，并说明相应奇延拓的周期。
\end{exercise}
\begin{exercise}{783}{常规·锯齿函数展开}
将 $f(x)=x$（$-\pi<x<\pi$）作 $2\pi$ 周期延拓，并令 $f(\pm\pi)=0$。利用奇性求傅里叶级数，写出它在周期端点处的和。
\end{exercise}
\begin{exercise}{784}{基础·周期端点平均}
设 $f$ 以 $2$ 为周期，在 $[-1,0)$ 为 $1-3x$，在 $[0,1)$ 为 $x^2+1$。不计算系数，求傅里叶级数在 $x=-7,0,1/2$ 处的和。
\end{exercise}
\begin{exercise}{785}{基础·余弦级数偶延拓}
设 $f(x)=x$（$0\leq x\leq1$），将其作偶延拓再作周期 $2$ 延拓。求对应余弦级数在 $x=-1/2,1,2$ 处的和，无需计算系数。
\end{exercise}
\begin{exercise}{786}{基础·正弦延拓端点}
设 $f(x)=x$（$0<x<1$），$b_n=2\int_0^1x\sin(n\pi x)\,\dd x$。求 $\sum_{n=1}^{\infty}b_n\sin(n\pi x)$ 在 $x=-1/2,0,1$ 处的和。
\end{exercise}
\begin{exercise}{787}{基础·取点求级数和}
已知 $x^2=\pi^2/3+4\sum_{n=1}^{\infty}(-1)^n\cos(nx)/n^2$ 在 $[-\pi,\pi]$ 成立。取 $x=\pi$ 求 $\sum_{n=1}^{\infty}n^{-2}$。
\end{exercise}
\begin{exercise}{788}{常规·余弦系数计算}
将 $f(x)=x$（$0\leq x\leq1$）展开成周期 $2$ 的余弦级数，采用 $a_0/2+\sum_{n=1}^{\infty}a_n\cos(n\pi x)$ 的约定，求 $a_0,a_n$。
\end{exercise}
\begin{exercise}{789}{基础·方波跳跃平均}
设周期 $2$ 的函数在 $(-1,0)$ 等于 $-1$，在 $(0,1)$ 等于 $1$，在整数处取零。求其傅里叶级数在 $x=0,1,1/2$ 处的和。
\end{exercise}
\begin{exercise}{790}{常规·平方函数傅里叶展开}
将 $f(x)=x^2$（$-\pi\leq x\leq\pi$）作 $2\pi$ 周期延拓。求傅里叶级数，并利用 $x=0$ 的取值求 $\sum_{n=1}^{\infty}(-1)^{n-1}/n^2$。
\end{exercise}
\subsubsection{常规综合练习}
\begin{exercise}{791}{基础·必要条件与极限}
已知 $\sum_{n=1}^{\infty}(u_n-3n/(n+1))$ 收敛。求 $\lim_{n\to\infty}u_n$ 及 $\lim_{n\to\infty}(u_n^2-2u_n+3)$。
\end{exercise}
\begin{exercise}{792}{常规·错位级数求和}
已知 $\sum_{k=1}^nu_k=n/(2n+1)$。求 $u_n$ 及 $\sum_{n=1}^{\infty}(u_n+u_{n+1})$ 的和。
\end{exercise}
\begin{exercise}{793}{常规·三重分母裂项}
用 $\frac1{n(n+1)(n+2)}=\frac12(\frac1{n(n+1)}-\frac1{(n+1)(n+2)})$ 求从 $n=1$ 开始的级数和。
\end{exercise}
\begin{exercise}{794}{基础·幂和代入}
已知 $1^3+\cdots+n^3=(n(n+1)/2)^2$。求 $\sum_{n=1}^{\infty}(1^3+2^3+\cdots+n^3)^{-1/2}$ 的和。
\end{exercise}
\begin{exercise}{795}{基础·正项级数部分和}
设 $u_n\geq0$。利用单调有界定理说明 $\sum u_n$ 收敛当且仅当部分和数列有上界。
\end{exercise}
\begin{exercise}{796}{基础·相邻倒数比较}
利用 $1/(n+1)^2\leq1/n-1/(n+1)$，证明 $\sum_{n=1}^{\infty}1/(n+1)^2$ 收敛。
\end{exercise}
\begin{exercise}{797}{常规·有理化与比较}
判定 $\sum_{n=1}^{\infty}(1/\sqrt n-1/\sqrt{n+1})$ 的敛散性并求和。
\end{exercise}
\begin{exercise}{798}{基础·积分定义的通项}
设 $u_n=\int_n^{n+1}e^{-x}\,\dd x$。求 $\sum_{n=1}^{\infty}u_n$ 的和。
\end{exercise}
\begin{exercise}{799}{基础·无穷小阶比较}
用等价无穷小分别判定 $\sum_{n=1}^{\infty}\sin(1/n)$ 与 $\sum_{n=1}^{\infty}(1-\cos(1/n))$ 的敛散性。
\end{exercise}
\begin{exercise}{800}{常规·极限指数比较}
用极限比较法判定 $\sum_{n=1}^{\infty}(n+2)^n/n^{n+2}$ 的敛散性。
\end{exercise}
\begin{exercise}{801}{基础·交错分母}
判定 $\sum_{n=1}^{\infty}(-1)^{n+1}/(\sqrt n+1)$ 的收敛类型。
\end{exercise}
\begin{exercise}{802}{基础·平移交错分母}
判定 $\sum_{n=1}^{\infty}(-1)^n/(\sqrt{2n}+1)$ 的收敛类型。
\end{exercise}
\begin{exercise}{803}{常规·整数相位化简}
利用 $\sin(n\pi+t)=(-1)^n\sin t$ 判定 $\sum_{n=1}^{\infty}\sin(n\pi+1/n)$ 的收敛类型。
\end{exercise}
\begin{exercise}{804}{基础·奇偶子级数}
分别判定 $\sum_{n=1}^{\infty}1/(2n-1)$、$\sum_{n=1}^{\infty}1/(2n)$ 与 $\sum_{n=1}^{\infty}(1/(2n-1)-1/(2n))$ 的敛散性。
\end{exercise}
\begin{exercise}{805}{常规·分式变量收敛域}
求 $\sum_{n=1}^{\infty}\frac1n((1-x)/(1+x))^n$ 的实收敛域，明确排除 $x=-1$ 并检查比值等于 $\pm1$ 的情况。
\end{exercise}
\begin{exercise}{806}{基础·平移对数和函数}
求 $\sum_{n=1}^{\infty}(-1)^{n-1}(x+1)^n/n$ 的收敛域及和函数，检查两个端点。
\end{exercise}
\begin{exercise}{807}{基础·平移调和系数}
求 $\sum_{n=1}^{\infty}(x-2)^n/n$ 的收敛域，分别判断两个端点。
\end{exercise}
\begin{exercise}{808}{常规·裂项数项求和}
利用部分分式求 $\sum_{n=2}^{\infty}1/(n^2-1)$ 的和。
\end{exercise}
\begin{exercise}{809}{常规·阶乘加权数项和}
利用 $n/(n+1)!=1/n!-1/(n+1)!$，求 $\sum_{n=1}^{\infty}n/(n+1)!$ 的和。
\end{exercise}
\begin{exercise}{810}{基础·偶数幂与对数}
求 $\sum_{n=1}^{\infty}(-1)^{n-1}x^{2n}/n$ 的收敛域及和函数，检查两个端点。
\end{exercise}
\begin{exercise}{811}{常规·反常面积与求和}
设 $D_n=\{x\geq e^2,0\leq y\leq1/(x(\ln x)^{n+1})\}$，$n\geq1$。求面积 $A_n$，再利用对数展开求 $\sum_{n=1}^{\infty}(-1)^{n-1}A_n$ 的和。
\end{exercise}
\begin{exercise}{812}{基础·两个对数展开}
将 $\ln((1+x)/(1-x))$ 展开成 $x$ 的幂级数，给出收敛域并检查两个端点。
\end{exercise}
\begin{exercise}{813}{基础·反正切复合展开}
将 $x\arctan(x/2)$ 展开成 $x$ 的幂级数，求收敛域并检查两个端点。
\end{exercise}
\begin{exercise}{814}{常规·乘积系数提取}
设 $f(x)=xe^{-x^2}\sin2x$。将两个因子展开到所需阶，求 $f^{(4)}(0)$。
\end{exercise}
\begin{exercise}{815}{常规·幂级数乘积}
设 $|x|<1$。求 $(\sum_{n=1}^{\infty}x^n)^2$ 展开式中 $x^n$（$n\geq2$）的系数，并求四阶原点导数。
\end{exercise}
\begin{exercise}{816}{基础·方波傅里叶和}
设 $f(x)=\operatorname{sgn}(\cos x)$，规定 $\operatorname{sgn}(0)=0$。不计算系数，求其傅里叶级数在 $x=-\pi/2,0,3\pi$ 处的和。
\end{exercise}
\begin{exercise}{817}{基础·余弦常数项变式}
将 $f(x)=1-x$（$0\leq x\leq1$）作偶延拓并以 $2$ 为周期延拓。按 $a_0/2+\sum_{n=1}^{\infty}a_n\cos(n\pi x)$ 的约定，求 $a_0$ 与 $a_1$。
\end{exercise}
\begin{exercise}{818}{常规·周期分段函数小综合}
设周期 $2$ 的函数在 $-1<x\leq0$ 时为 $x^2$，在 $0<x\leq1$ 时为 $x-1$。求其傅里叶级数在 $x=0,1,1/2$ 处的和，并求常数项 $a_0/2$，其中 $a_0=\int_{-1}^1f(x)\,\dd x$。
\end{exercise}

\end{document}
